\documentclass[a4paper, 12pt]{amsart}
\usepackage[T1]{fontenc}
\usepackage{amsfonts,enumerate,array,amsmath,xcolor}
\usepackage{graphicx,float}
\usepackage[english]{babel}
\usepackage{fancyhdr}
\usepackage{tikz}
\usepackage{tikz-cd}


%Margins

\setlength{\textwidth}{\paperwidth}
\addtolength{\textwidth}{-2in}
\calclayout

%Theorem environments
\theoremstyle{definition}
\newtheorem{defn0}{Definition}[section]

\theoremstyle{theorem}
\newtheorem{prop0}[defn0]{Proposition}
\newtheorem{thm0}[defn0]{Theorem}
\newtheorem{lemma0}[defn0]{Lemma}
\newtheorem{corollary0}[defn0]{Corollary}
\newtheorem{example0}[defn0]{Example}
\newtheorem{remark0}[defn0]{Remark}
\newtheorem{conjecture0}[defn0]{Conjecture}

\newenvironment{definition}{ \begin{defn0}}{\end{defn0}}
\newenvironment{proposition}{\bigskip \begin{prop0}}{\end{prop0}}
\newenvironment{theorem}{\bigskip \begin{thm0}}{\end{thm0}}
\newenvironment{lemma}{\bigskip \begin{lemma0}}{\end{lemma0}}
\newenvironment{corollary}{\bigskip \begin{corollary0}}{\end{corollary0}}
\newenvironment{example}{ \begin{example0}\rm}{\end{example0}}
\newenvironment{remark}{ \begin{remark0}\rm}{\end{remark0}}
\newenvironment{conjecture}{\begin{conjecture0}}{\end{conjecture0}}

\newcommand{\defref}[1]{Definition~\ref{#1}}
\newcommand{\propref}[1]{Proposition~\ref{#1}}
\newcommand{\thmref}[1]{Theorem~\ref{#1}}
\newcommand{\lemref}[1]{Lemma~\ref{#1}}
\newcommand{\corref}[1]{Corollary~\ref{#1}}
\newcommand{\exref}[1]{Example~\ref{#1}}
\newcommand{\secref}[1]{Section~\ref{#1}}
\newcommand{\remref}[1]{Remark~\ref{#1}}
\newcommand{\conjref}[1]{Conjecture~\ref{#1}}

%Commands

\newcommand{\Q}{\mathbb{Q}}
\newcommand{\Z}{\mathbb{Z}}
\newcommand{\Zmod}[1]{\mathbb{Z}/#1\mathbb{Z}}
\renewcommand{\P}{\mathfrak{P}}
\newcommand{\Div}{\text{Div}}
\newcommand{\GF}[1]{\mathbb{F}_{#1}}
\newcommand{\Jac}{\text{Jac}}
\DeclareFontFamily{U}{wncy}{}
\DeclareFontShape{U}{wncy}{m}{n}{<->wncyr10}{}
\DeclareSymbolFont{mcy}{U}{wncy}{m}{n}
\DeclareMathSymbol{\Sh}{\mathord}{mcy}{"58}
\newcommand{\Gal}{\text{Gal}}
\newcommand{\Ker}{\text{Ker}}
\renewcommand{\Im}{\text{Im}}
\newcommand{\Hom}{\text{Hom}}


%Fixing bibliography
%\DeclareUnicodeCharacter{202F}{FIX ME!!!!}

%%%%%%%%%%%%%%%%%%%%%%%%%%%%%%%%%%%%%%%%%%%%%%%%%%%%%%%%%%%%%
%%%%%%%%%%%%%%%%%%%%%%%%%%%%%%%%%%%%%%%%%%%%%%%%%%%%%%%%%%%%%


\begin{document}
	%\frontmatter
	\thispagestyle{empty}
	\pagenumbering{roman} \setcounter{page}{0}
	\begin{figure}[ht]
		\centering\hbox{
			\includegraphics[width=0.6\textwidth]{marca_pos_rgb.png}\hspace{2cm}
			\includegraphics[width=0.4\textwidth]{Logo_uab.png}}
	\end{figure}
	\vglue 2truecm
	\centerline{\bf \large ADVANCED MATHEMATICS}
	\vglue 0.1truecm
	\centerline{\bf  MASTER'S FINAL PROJECT}
	\vglue 2truecm
	\hrule
	\vglue 1truecm
	\centerline{\bf \Large ON THE MORDELL - WEIL RANK}
	\vspace{0.4cm}
	\centerline{\bf \Large OF HYPERELLIPTIC CURVES}
	\vglue 1truecm
	\hrule
	\vglue 1truecm
	\noindent {\it Author:} \hfill {\it Supervisor:}
	\newline
	\noindent {\color{red}{Guillem Mata Carmona}} \hfill {\color{red}{Marc Masdeu Sabaté}}
	\vglue 3truecm
	\centerline{\bf \Large Facultat de Matemàtiques i Informàtica}
	\vglue 2truecm \hfill{\today}
	
	\newpage
	
	\tableofcontents
	
	\newpage
	\thispagestyle{empty}
	\section*{Abstract}
	Lorem ipsum
	
	
	
	\newpage
	\pagenumbering{arabic}
	\pagestyle{fancy}
	
	\section{Introduction}
	
	%The following is just a copy of a draft
	Let us explain in detail the method described in Silverman's book \cite{Silverman_2009} to compute the rank of an elliptic curve. Let $E/K$ be an elliptic curve defined over some number field $K$. By the Mordell-Weil theorem, the group of $K$-rational points of $E$, $E(K)$, is finitely generated.
	
	%The following is also a copy of a (longer) draft
	
	We begin by recalling the following classical situation: let $K$ be a number field, and $E/K$ some elliptic curve defined over $K$. The set of $K$-rational points, denoted by $E(K)$ has a structure of abelian group, and moreover we have:
	\begin{theorem}[Mordell-Weil]
		The group of $K$-rational points of an elliptic curve is finitely generated. Therefore, we have an isomorphism:
		\begin{equation*}
			E(K)\cong \mathbb Z^r\oplus E(K)_{tors},
		\end{equation*}
		where $r\in\mathbb{Z}_{\ge 0}$ is called the rank of $E$ and $E(K)_{tors}$ is the torsion subgroup of $E(K)$, which is a finite abelian group.
	\end{theorem}
	It is well known that in order to prove the Mordell-Weil theorem, one can use the theory of height functions to reduce to proving the fact that the group $E(K)/mE(K)$ is finite for $m\ge 2$, which is known as the weak Mordell-Weil theorem. We describe in this section the method for computing the rank of an elliptic curve explained in J. Silverman's book. Essentially, the method consists on proving the weak Mordell-Weil theorem by embedding the group $E(K)/mE(K)$ into a certain (finite) cohomology group which allows explicit computations to be done.
	
	From a theoretical point of view, it is possible (and necessary) to make the assumption that $E[m]\subseteq E(K)$ without loss of generality when proving the WMW theorem, due to the following fact: if $L/K$ is a finite extension and we know that $E(L)/mE(L)$ is finite, then $E(K)/mE(K)$ is finite. If we instead want to compute explicitly the rank of $E$, we have to take the assumption $(*)$ in order for our method to work, but if sdñlkjañd
	
	Thus we make the assumption that $E[m]\subseteq E(\mathbb{Q})$. Let us recall some basic constructions for elliptic curves.
	
	Fix once and for all an algebraic closure of $K$, $\overline{K}$. We denote the absolute Galois group of $K$ by $G_{\overline{K}/K}:=\text{Gal}(\overline{K}/K)$, which acts on $Q\in E(\overline{K})$ by applying $\sigma\in G_{\overline{K}/K}$ to the coordinates of $Q$. Given $P\in E(K)$ and $\sigma\in G_{\overline{K}/K}$, we can find $Q\in E(\overline{K})$ such that $mQ=P$, and thus we obtain a pairing, called the Kummer pairing:
	\begin{align*}
		\kappa:E(K)\times G_{\overline{K}/K}&\longrightarrow E[m]\\
		(P,\sigma)&\mapsto \kappa(P,\sigma)=Q^{\sigma}-Q
	\end{align*}
	which is bilinear and its left kernel is $mE(K)$. Therefore we obtain a homomorphism:
	\begin{align*}
		\delta_E:E(K)/mE(K)&\hookrightarrow\text{Hom}(G_{\overline{K}/K},E[m])\\
		[P]&\mapsto \delta_E(P)=\kappa(P,\cdot\ )
	\end{align*}
	We also have the Weil pairing $e_m:E[m]\times E[m]\rightarrow \mu_m$ defined as follows: let $T\in E[m]$, and take a function $f\in \overline{K}(E)$ such that:
	\begin{equation*}
		\text{div}(f)=m(T)-m(O).
	\end{equation*}
	
	
	
	asdf
	
	There is a bilinear pairing
	\begin{equation*}
		b:E(K)/mE(K)\times E[m]\rightarrow K^*/(K^*)^m,
	\end{equation*}
	such that
	\begin{equation*}
		e_m(\delta_E(P),T)=\delta_K(b(P,T))
	\end{equation*}
	as elements in $\text{Hom}(G_{\overline{K}/K},\mu_m)$, so that for every $P,T\in E[m]$ and every $\sigma\in G_{\overline{K}/K}$ we have $e_m(\delta_E(P)(\sigma),T)=\delta_K(b(P,T))(\sigma)$. The following properties of the pairing $b$ will allow us to prove the WMW theorem explicitly:
	\begin{theorem}
		\phantom{a}
		
		\begin{itemize}
			\item The pairing $b$ is non degenerate on the left.
			\item $K(S,m)$
			\item computation
		\end{itemize}
	\end{theorem}
	The fact that $b$ is non degenerate on the left means that there is an injective homomorphism
	\begin{equation*}
		E(K)/mE(K)\hookrightarrow \text{Hom}(E[m],K^*/(K^*)^m).
	\end{equation*}
	This is the kind of embedding that we would like to use in order to deduce that $E(K)/mE(K)$ is finite. However, $K^*/(K^*)^m$ is not finite so we cannot conclude. However, by part $b)$ of the previous theorem we find that we actually have an injective homomorphism
	\begin{equation*}
		E(K)/mE(K)\hookrightarrow \text{Hom}(E[m],K(S,m)),
	\end{equation*}
	and since both $E[m]$ and $K(S,m)$ are finite groups, we conclude that $E(K)/mE(K)$ is finite. Since we have an explicit way of computing $b$, we can try to determine the image of $E(K)/mE(K)$ in $\text{Hom}(E[m],K(S,m))$ where, as we will see, computations are easier.
	
	First of all, we know that
	\begin{equation*}
		E[m]\cong \mathbb{Z}/m\mathbb{Z}\times\mathbb{Z}/m\mathbb{Z},
	\end{equation*}
	so we begin by choosing a pair of generators $T_1$, $T_2$ of $E[m]$. An element $\varphi\in \text{Hom}(E[m],K(S,m))$ is determined by its values at $T_1$ and $T_2$, which is a pair of elements in $K(S,m)$ so we have an isomorphism:
	\begin{align*}
		\text{Hom}(E[m],K(S,m))&\rightarrow K(S,m)\times K(S,m)\\
		\varphi&\mapsto(\varphi(T_1),\varphi(T_2))
	\end{align*}
	The conclusion of this discussion is that we have a well defined injective homomorphism
	\begin{align*}
		\psi:E(K)/mE(K)&\hookrightarrow K(S,m)\times K(S,m)\\
		[P]&\mapsto(b(P,T_1),b(P,T_2))
	\end{align*}
	With these identifications, we want to check for each of the (finitely many) pairs $(b_1,b_2)\in K(S,m)\times K(S,m)$, whether they come from an element in $E(K)/mE(K)$, i.e., if they lie in the image of $\psi$. This will happen if and only if there exists some $P\in E(K)$ such that $b_1=b(P,T_1)$, $b_2=b(P,T_2)$. Since we know how to compute $b(P,T)$ up to $m$-th powers, this is equivalent to the existence of $P\in E(K)$ and $z_1,z_2\in K^*$ such that
	\begin{equation*}
		b_1z_1^m=f(P,T_1),\ b_2z_2^m=f(P,T_2).
	\end{equation*}
	\subsection{Example}
	We use the method described above to compute the group of rational points of the curve:
	\begin{equation*}
		E:y^2=x^3-21x-20=(x+1)(x-5)(x+4)
	\end{equation*}
	First of all, we note that the 2-torsion is defined over $\mathbb{Q}$, since:
	\begin{equation*}
		E[2]\cong \mathbb{Z}/2\mathbb{Z}\oplus\mathbb{Z}/2\mathbb{Z}=\{\infty,(-1,0),(5,0),(-4,0)\}
	\end{equation*}
	The discriminant of this curve is $\Delta=419904=2^6 3^8$, so the curve has bad reduction at $2$ and $3$. A quick computation shows that 
	
	\newpage
	\section{Jacobians of curves}
	
	\newpage
	\section{The Mordell-Weil theorem}
	
	In this chapter, we are going to prove the Mordell-Weil theorem for abelian varieties over number fields, which will be crucial for our later work. The theorem states the following:
	
	\begin{theorem}\label{mordellweiltheorem}(Mordell-Weil)
		
		Let $A$ be an abelian variety defined over a number field $K$. Then the group of $K$-rational points $A(K)$ is finitely generated. In particular, there exists some integer $r\ge 0$ called the Mordell-Weil rank of $A$ over $K$, and integers $n_1,\ldots,n_s\ge 2$ such that:
		\begin{equation*}
			A(K)\cong \Z^r\oplus \Zmod{n_1}\oplus\ldots\oplus\Zmod{n_s}
		\end{equation*}
	\end{theorem}
	In the next chapters we will turn into the problem of computing the Mordell-Weil rank. For now, let us make the following observation: if $A$ is a finitely generated abelian group, then it is easy to see that the quotient group $A/mA$, where $mA$ denotes the image of $A$ under the multiplication-by-$m$ map, is finite. The converse is not true: take the multiplicative abelian group $\mathbb{R}^*$, which is not finitely generated. However, we have:
	\begin{equation*}
		\mathbb{R}^*/(\mathbb{R}^*)^2\cong\{\pm 1\},
	\end{equation*}
	which is finite. In the particular case of the group $A(K)$, the interplay between the geometry of the variety and the arithmetic of its points allows us to construct certain functions that prove the converse. For this reason, the proof of the Mordell-Weil theorem has two steps: in the first section of this chapter, we will show that the quotient $A(K)/mA(K)$ is finite, by using some results from algebraic number theory. In the second section of this chapter, we will explain the construction of the canonical height functions that encode the structure of $A(K)$ and will use them to deduce \thmref{mordellweiltheorem} from the finiteness of the quotient $A(K)/mA(K)$.
	
	\subsection{The Weak Mordell-Weil theorem}
	
	As we claimed, in this section we will prove the following key theorem for the proof of \thmref{mordellweiltheorem}:
	\begin{theorem}\label{weakmordellweil}(Weak Mordell-Weil)
		
		Let $A$ be an abelian variety defined over a number field $K$. Let $m\ge 2$ be a positive integer, and denote by $mA(K)$ the image of $A(K)$ by the multiplication-by-$m$ morphism. Then the group
		\begin{equation*}
			A(K)/mA(K)
		\end{equation*}
		is finite.
	\end{theorem}
	
	In all the following, we will consider $m$ as being fixed and will call $A(K)/mA(K)$ the weak Mordell-Weil group of $A$ over $K$.
	
	The strategy for the proof of \thmref{weakmordellweil} is to use the structure of $A$ to construct an injective map from $A(K)/mA(K)$ to a set, which will turn out to be finite after some additional work.
	
	First of all, we will explain why it suffices to show \thmref{weakmordellweil} for a finite extension of $K$.
	
	\begin{lemma}\label{finitereductionlemma}
		Let $K'/K$ be a finite Galois extension of $K$. Assume that the group
		\begin{equation*}
			A(K')/mA(K')
		\end{equation*}
		is finite. Then the group $A(K)/mA(K)$ is also finite.
	\end{lemma}
	\begin{proof}
		Let $G$ be the Galois group of $K'$ over $K$. The inclusion $A(K)\subset A(K')$ induces a map $A(K)\rightarrow A(K')/mA(K')$. Note that points in $mA(K)$ go to $mA(K')$, so we have a natural map
		\begin{equation*}
			A(K)/mA(K)\rightarrow A(K')/mA(K').
		\end{equation*}
	Let us compute its kernel: if the class modulo $mA(K)$ of $x\in A(K)$ is sent to zero by this map, then it means that $x$ is in $mA(K')$. Therefore, the kernel is contained in the quotient $A(K)\cap mA(K')/mA(K)$. The reverse inclusion is clear, so we find that the kernel is equal to:
	\begin{equation*}
		B:=A(K)\cap mA(K')/mA(K).
	\end{equation*}
	We want to show that this kernel is finite. To do so choose, for any class $b\in B$, a point $x\in A(K)\cap mA(K')$ such that $[x]=b$. Then choose a point $y\in A(K')$ such that $[m]y=x$. For any $\sigma\in \Gal(K'/K)$, define $\kappa_x(\sigma)=y^\sigma-y$. Observe that:
	\begin{equation*}
		[m]\kappa_x(\sigma)=([m]y)^\sigma-[m]y=x^\sigma-x=x-x=0,
	\end{equation*}
	so $\kappa_x(\sigma)\in A[m]$. We can thus define a map, for each $x\in A(K)$:
	\begin{align*}
		\kappa_x:\Gal(K'/K)&\rightarrow A[m]\\
		\sigma\mapsto \kappa_x(\sigma).
	\end{align*}
	These type of maps will also be important later. For now, we do not care about the dependency of their value on the choices of $x$ and $y$. We will be concerned on this problem and the properties of the $\kappa_x$ after this proof.
	
	Note that the choices of $y$ and $x$ for each $b$ yield a map:
	\begin{equation}\label{homset}
		B\rightarrow \Hom_{\text{set}}(\Gal(K'/K),A[m]),
	\end{equation}
	which is not in principle a morphism of groups, just a set map. What we can actually show is that this assignation is injective: suppose $b,b'$ are classes in $B$ such that $k_x$ and $k_{x'}$ are equal, where $x$ is our chosen representative for $b$ and $x'$ is our chosen representative for $b'$. Then, if $y,y'$ are the chosen points in $A(K')$ such that $[m]y=x$ and $[m]y'=x'$, we obtain the equality $y^\sigma-y=(y')^\sigma-y'$, so:
	\begin{equation*}
		(y-y')^\sigma=y-y',
	\end{equation*}
	and hence we conclude that $y-y'$ is in $A(K)$. But then, observe that $x-x'=[m]y-[m]y'=[m](y-y')$, so $x-x'\in mA(K)$, so $b=b'$ as classes of $B$, which is what we wanted.
	
	Since the extension $K'/K$ is finite and the group $A[m]$ is finite (isomorphic to $(\Zmod{m})^{2g}$ where $g$ is the dimension of $A$), we conclude that $B$ is finite from the existence of the map \eqref{homset}.
	
	The group $A(K)/mA(K)$ fits in the exact sequence:
	\begin{equation*}
		\begin{tikzcd}
			0\arrow{r}
			& B\arrow{r}
			& A(K)/mA(K)\arrow{r}
			& A(K')/mA(K'),
		\end{tikzcd}
	\end{equation*}
	so we find that $A(K)/mA(K)$ contains a finite subgroup of finite index, so it must be finite.
	\end{proof}
	Using this lemma we can assume, extending $K$ by a finite extension if needed, that $A[m]$ is contained in $A(K)$. By the same argument we may also assume that the $m^\text{th}$ roots of unity are in $K$. Actually the assumption on the $m^\text{th}$ roots of unity is automatically satisfied for fields such that $A[m]$ is contained in $A(K)$, which is usually proven using a tool called the Weil pairing. We will discuss the Weil pairing for Jacobian varieties and its properties in the next chapter, so at this point we take these assumptions as if they were independent. In the following, we will refer to these assumptions as $(*)$.
	
	Under these assumptions we can define a map with the properties that we want. Fix once and for all an algebraic closure $\overline{K}$ of $K$ and denote by $G_K$ the absolute Galois group $\Gal(\overline{K}/K)$.
	\begin{definition}
		Let $x$ be a point in $A(K)$. Choose some $y\in A(\overline{K})$ such that $[m]y=x$. For any $\sigma\in G_K$, define $\kappa(x,\sigma)=y^\sigma-y$. The Kummer pairing is the map:
		\begin{align*}
			\kappa:A(K)\times G_K&\rightarrow A[m]\\
			(x,\sigma)&\mapsto \kappa(x,\sigma)
		\end{align*}
	\end{definition}
	We have to check that the Kummer pairing is well defined, i.e., it does not depend on the choice of $y\in A(\overline{K})$ and the image lies in $A[m]$. First, assume that $y,y'$ are two points in $A(\overline{K})$ such that $[m]y=x=[m]y'$. Then, we have:
	\begin{equation*}
		0=[m]y-[m]y'=[m](y-y'),
	\end{equation*}
	so $y-y'$ is a point of $m$ torsion. Since we assumed that $A[m]\subset A(K)$, we have, for any $\sigma\in G_K$:
	\begin{align*}
		(y^\sigma-y)-((y')^\sigma-y')=(y-y')^\sigma - (y-y')=(y-y')-(y-y')=0,
	\end{align*}
	because $G_K$ fixes $A(K)$, so $\kappa(x,\sigma)$ does not depend on the choice of $y$. Next, we show that the image lies in the $m$ torsion points of $A$. Indeed, choose any $y\in A(\overline{K})$ such that $[m]y=x$, then:
	\begin{align*}
		[m]\kappa(x,\sigma)&=[m](y^\sigma-y)=[m](y^\sigma)-[m]y\\
		&=([m]y)^\sigma-[m]y=x^\sigma-x=x-x=0,
	\end{align*}
	as we wanted. Now that we have seen the well-definedness of $\kappa$, let us explore some of its properties.
	\begin{proposition}\label{kummerproperties}
		Assume $(*)$. Then the Kummer pairing is a bilinear map, whose left kernel is $mA(K)$. Let $L$ be the extension of $K$ obtained by adjoining the coordinates of all elements $y\in A(\overline{K})$ such that $[m]y=x$, for some $x\in A(K)$. Then the right kernel of this map is the group $\Gal(\overline{K}/L)$.
	\end{proposition}
	\begin{proof}
		Let us show that the pairing is bilinear. Let $x,x_1,x_2$ be in $A(K)$, and let $\sigma,\tau$ be in $G_K$. Choose $y,y_1,y_2$ in $A(\overline{K})$ such that $[m]y=x,[m]y_1=x_1$ and $[m]y_2=x_2$. Then:
		\begin{align*}
			\kappa(x,\sigma\tau)&=y^{\sigma\tau}-y=y^{\sigma\tau}-y^\tau+y^\tau-y\\
			&=(y^\sigma-y)^\tau + y^\tau-y\\
			&=\kappa(x,\sigma)^\tau + \kappa(x,\tau)\\
			&=\kappa(x,\sigma) + \kappa(x,\tau),
		\end{align*}
		where the last equality follows from the fact that $\kappa(x,\sigma)$ is a point of $m$-torsion, which is contained in $A(K)$ This shows linearity on the left. For linearity on the right, we compute:
		\begin{align*}
			\kappa(x_1,\sigma)+\kappa(x_2,\sigma)&=(y_1^\sigma-y_1)+(y_2^\sigma-y_2)\\
			&=(y_1+y_2)^\sigma - (y_1+y_2)\\
			&=\kappa(x_1+x_2,\sigma),
		\end{align*}
		because $[m](y_1+y_2)=x_1+x_2$. This shows linearity on the right, and we are done.
		
		Now we compute the left kernel: suppose that $x\in mA(K)$, and choose $y\in A(K)$ such that $[m]y=x$. Then we have, for all $\sigma\in G_K$:
		\begin{align*}
			\kappa(x,\sigma)=y^\sigma-y=y-y=0,
		\end{align*}
		which shows that $mA(K)$ is contained in the left kernel. To show the reverse inclusion, assume that $x\in A(K)$ is such that $\kappa(x,\sigma)=0$ for all $\sigma\in G_K$. Choose some $y\in A(\overline{K})$ such that $[m]y=x$, since $\kappa(x,\sigma)=0$, we have that $y^\sigma=y$ for all $\sigma\in G_K$. Therefore, $y$ is in $A(K)$ and thus $x=[m]y$ is in $mA(K)$, as we wanted.
		
		Finally, we compute the right kernel: let $x$ be any point in $A(K)$ and assume that $\sigma\in\Gal(\overline{K}/L)$. Choose some $y\in A(\overline{K})$ such that $[m]y=x$, then we have:
		\begin{equation*}
			\kappa(x,\sigma)=y^\sigma-y=y-y=0,
		\end{equation*}
		because by definition the coordinates of $y$ are in $L$, so $\sigma$ fixes it. Now let $\sigma$ be a Galois morphism such that $\kappa(x,\sigma)=0$, for all $x\in A(K)$. To show that $\sigma$ is in $\Gal(\overline{K}/L)$, we must show that it fixes every element in $L$. Let $y\in A(\overline{K})$ be a point such that $[m]y\in A(K)$, say it equals some $x\in A(K)$. Then $\kappa(x,\sigma)=0$, so $y^\sigma=y$, as we wanted. Since $\sigma$ fixes the generators of $L$, it fixes all of $L$.
	\end{proof}
	\begin{remark}
		\propref{kummerproperties} implies that there is a nondegenerate bilinear pairing
		\begin{equation*}
			A(K)/mA(K)\times\Gal(L/K)\rightarrow A[m].
		\end{equation*}
		We then obtain a morphism
		\begin{equation*}
			A(K)/mA(K)\hookrightarrow\Hom(\Gal(L/K),A[m])
		\end{equation*}
		given by fixing the first component of the pairing, which is injective by the nondegeneracy of the pairing. Since $A[m]$ is a finite group, if we can show that the extension $L/K$ described in \propref{kummerproperties} is finite, we will be done. That is exactly what we are going to do for the remain of this section.
	\end{remark}
	\begin{remark}
		If we drop our assumptions on the $m$-torsion and the $m^\text{th}$ roots of unity, then the arguments given in \lemref{finitereductionlemma} and \propref{kummerproperties} do not work anymore. Instead, what we obtain are well defined elements in cohomology. We will delve into this situation in the next chapter, where we will translate the arguments given here to the language of Galois cohomology.
	\end{remark}
	In order to show that the extension $L/K$ is finite, we must study some of its properties. Note that $L$ is the compositum of all the extension of the form $K(y)$, where $[m]y=x$ for some $x\in A(K)$, so we first study these extensions. To simplify notation, we will denote by $A(y)$ the group of $K(y)$-rational points of $A$.
	\begin{lemma}
		With the same notation and assumptions as above, the extension $K(y)$ is Galois over $K$, with Galois group isomorphic to a subgroup of $A[m]$.
	\end{lemma}
	\begin{proof}
		We want to show that (the coordinates of) all the Galois conjugates of $y$ already lie in $K(y)$. Let $\sigma$ be an element of $G_K$. Then we see that $y^\sigma=\kappa(x,\sigma)+y$, by definition. We have already seen that $\kappa(x,\sigma)$ is an element of $m$-torsion, which by assumption lies in $A(K)$, so any Galois conjugate of $y$ differs by an element already in $A(K)$, and thus it lies in $A(y)$. Thus we have that the coordinates of $y^\sigma$ lie in $K(y)$, and so the extension is Galois.
		
		Moreover, we have a well defined map
		\begin{align*}
			\Gal(K(y)/K)&\rightarrow A[m]\\
			\sigma &\mapsto y^\sigma-y
		\end{align*}
		which is a group morphism by \propref{kummerproperties}. Observe that any two morphisms $\sigma,\tau$ which map to the same element satisfy $y^\sigma-y=y^\tau-y$, so $y^\sigma=y^\tau$. Since $y$ generates the extension $K(y)$, we have that $\sigma,\tau$ must be equal, so this map is injective. Therefore $\Gal(K(y)/K)$ is isomorphic to a subgroup of $A[m]$, as we wanted.
	\end{proof}
	Note that we have proved that the extensions $K(y)/K$ are all finite abelian Galois extensions of exponent dividing $m$.
	
	We can also study the ramification of such extensions, but to do that we will need a result on the $m$-torsion of an abelian variety.
	
	\begin{theorem}\label{torsioninjects}
		Let $A$ be an abelian variety defined over a number field $K$. Let $\P$ be a non-archimedean prime of $K$ at which $A$ has good reduction. Let $\tilde{k}$ be the residue field of $\P$, with characteristic $p$. Let $m\ge 1$ be an integer coprime with $p$. Then the reduction map
		\begin{equation*}
			A(K)[m]\rightarrow \tilde{A}(\tilde{k})
		\end{equation*}
		is injective.
	\end{theorem}
	\begin{proof}
		The proof of this theorem involves the theory of formal groups, which we will not explain in this thesis for brevity. For a complete proof, see \cite{Hindry_Silverman_2000}.
	\end{proof}
	
	With this important fact, we can say some things about the ramification of $K(y)$.
	\begin{proposition}\label{ramifS}
		Let $m\ge 1$ be an integer, and let $S$ be the (finite) set of primes in $K$ consisting of the primes where $A$ has bad reduction and the primes dividing $m$. Then the extensions $K(y)/K$ described above are unramified outside $S$.
	\end{proposition}
	\begin{proof}
		Let $\P$ be any prime of $K$ not in $S$. In particular, $A(K)$ has good reduction at $\P$. Let $\P'$ be an extension of $\P$ to $K(y)$. Then $A(y)$ has good reduction at $\P'$, so we can consider the reduction map
		\begin{equation*}
			A(y)\rightarrow \tilde{A}_{\P'}(\tilde{k}_{\P'}).
		\end{equation*}
		To show that $K(y)$ is unramified at $\P$, we must show that $\P'$ has trivial inertia subgroup. Any morphism $\sigma\in\Gal(K(y)/K)$ in the decomposition subgroup for $\P$ descends to a morphism $\tilde{\sigma}$ for the residue field of $\P'$ over the residue field of $\P$. In this sense, a morphism $\sigma$ belongs to the inertia subgroup if and only if it reduces to the identity morphism. Let $\sigma$ be an element of the inertia subgroup for $\P'$. Since $\P$ does not divide $m$, we observe that the characteristic of the residue field of $\P$ is coprime to $m$. Using the definition of the inertia subgroup, we can compute:
		\begin{equation*}
			\widetilde{y^\sigma-y}=\tilde{y}^{\tilde{\sigma}}-\tilde{y}=\tilde{y}-\tilde{y}=\tilde{0},
		\end{equation*}
		since $\sigma$ acts trivially modulo $\P$. But now, our assumptions on $\P$ imply that the reduction map
		\begin{equation*}
			A(y)[m]\rightarrow \tilde{A}_{\P'}(\tilde{k}_{\P'})
		\end{equation*}
		is injective, by \thmref{torsioninjects}. By our assumption $(*)$, we have that $y^\sigma-y$ is a point of $m$-torsion. Therefore, we find that $y^\sigma-y=0$, so $y^\sigma=y$ and thus $\sigma$ is the identity, because $y$ generates $K(y)$. This shows that the inertia subgroup for $\P'$ is trivial and, since $\P$ was arbitrary, we are done.
	\end{proof}
	Note that we automatically have shown that $L$ is unramified outside $S$, because $L$ is the compositum of the extensions $K(y)$.
	
	The final step in the proof of \thmref{weakmordellweil} is to prove a finiteness result on extensions satisfying the properties of $K(y)/K$ that we have discussed. From Kummer theory, we know that abelian extensions $K'$ of $K$ of exponent $m$ correspond to subgroups $B$ of $K^*$ containing $(K^*)^m$, and this correspondence is given by
	\begin{equation*}
		B\mapsto K(\sqrt[m]{B}),\quad\text{and } K'\mapsto (K'^*)^m\cap K^*,
	\end{equation*}
	so we shall study extensions of the form $K(\sqrt[m]{a})$.
	
	\begin{lemma}
		Under the assumption $(*)$, let $\alpha\in K^*$. Let $K'=K(\sqrt[m]{\alpha})$ and let $\P$ be a prime of $K$ not dividing $m$. Then $K'/K$ is unramified at $\P$ if and only if the exact power of $\P$ dividing $\alpha$ is a multiple of $m$.
	\end{lemma}
	\begin{proof}
		Let $\omega:=\sqrt[m]{\alpha}$. The minimal polynomial for $\omega$ over $K$ is $f(x):=x^m-\alpha$, so we have that $f'(\omega)=m\omega^{m-1}$. The Galois conjugates of an element in $K'$ over $K$ correspond to multiplication by the $m^\text{th}$ roots of unity, so the Galois conjugates of $f'(\omega)$ are:
		\begin{equation*}
			m\omega^{m-1}, m\zeta_m\omega^{m-1},\ldots,m\zeta_m^{m-1}\omega^{m-1}.
		\end{equation*}
		Since the $m^\text{th}$ roots of unity are the roots of the polynomial $x^m-1$, the absolute value of the product of all the $m^\text{th}$ roots of unity is 1. Therefore, we have that:
		\begin{equation*}
			N_{K'/K}(f'(\omega))=\prod_{i=0}^{m-1}m\omega^{m-1}\zeta_m^i=\pm m^m\omega^{m(m-1)}=\pm m^m\alpha^{m-1},
		\end{equation*}
		so the discriminant of the order $\mathcal{O}_K[\omega]$ is $\pm m^m\alpha^{m-1}$. Hence, the discriminant of $K'/K$ divides $\pm m^m\alpha^{m-1}$ (it may not be equal because the order $\mathcal{O}_K[\omega]$ may not be maximal) and, since $\P$ does not divide $m$ by assumption, if $\P$ does not divide $\alpha$ (so the exact power of $\P$ dividing $\alpha$ is zero) then $K'/K$ is unramified at $\P$. For this reason, it suffices to show the case where $\P$ divides $\alpha$.
		
		Assume that the power of $\P$ dividing $\alpha$ is a multiple of $m$, say $mt$. Find an element $\pi\in K$ such that $\text{ord}_\P(\pi)=1$. Then $\P$ does not divide the element $\beta:=\pi^{-mt}\alpha$, but
		\begin{equation*}
			\beta=(\pi^{-t}\omega)^m,
		\end{equation*}
		so $K(\sqrt[m]{\alpha})=K(\sqrt[m]{\beta})$. Therefore, we are in the same situation as before and $K'/K$ is unramified at $\P$.
		
		Conversely, assume that the power of $\P$ dividing $\alpha$ is not a multiple of $m$, and call it $r$. We can then write the ideal generated by $\alpha$ as $\alpha\mathcal{O}_K=\P^r\mathfrak{Q}$, where $\mathfrak{Q}$ is coprime to $\P$. Then if we write $\P$ in $K'$ as a product of prime ideals, say $\P\mathcal{O}_{K'}=\P_1^{e_1}\cdot\ldots\cdot\P_s^{e_s}$, we obtain:
		\begin{equation*}
			\alpha\mathcal{O}_{K'}=\P_1^{re_1}\cdot\ldots\cdot\P_s^{re_s}\cdot\mathfrak{Q}',
		\end{equation*}
		where $\mathfrak{Q}'$ is some prime ideal in $K'$ coprime to the $\P_i$. Since $\alpha=\omega^m$ in $K'$, we have that $\alpha\mathcal{O}_{K'}=(\omega\mathcal{O}_{K'})^m$, so $m$ must divide every $re_i$. Since $m$ does not divide $r$ by assumption, it must divide each $e_i$, so each of the $e_i$ is strictly larger than 1. Thus $K'/K$ is ramified at $\P$, as we wanted.
	\end{proof}
	Now we recall a couple fundamental results from algebraic number theory that we will use to finish the proof. Let $K$ be a number field. Then the classes of ideals of $\mathcal{O}_K$ modulo principal ideals form an abelian group.
	\begin{theorem}
		The group of ideal classes of $K$ is finite.
	\end{theorem}
	Now let $S$ be a finite set of nonarchimedean primes of $K$. Define the ring of $S$-integers as:
	\begin{equation*}
		\mathcal{O}_{K,S}:=\{x\in K: v_\P(x)\ge 0\text{ for all non-archimedean } \P\not\in S\}.
	\end{equation*}
	Similarly, define the group of $S$-units as:
	\begin{equation*}
		\mathcal{O}_{K,S}^*:=\{x\in K^*: v_\P(x)=0\text{ for all non-archimedean } \P\not\in S\}.
	\end{equation*}
	Then we have:
	\begin{theorem}(Dirichlet's unit theorem)
		Let $r_1$ be the number of real embeddings of $K$ and $r_2$ the number of pairs of complex embeddings of $K$. Then $\mathcal{O}_{K,S}^*$ is a finitely generated abelian group of rank $r_1+r_2-1+|S|$, and the torsion subgroup corresponds to the roots of unity in $K$.
	\end{theorem}
	\begin{theorem}\label{maxabextension}
		Let $K$ be a number field, $m$ an integer and $S$ a finite set of non-archimedean primes of $K$. Assume that $K$ contains a primitive $m^{\text{th}}$ root of unity, and that $S$ contains all primes of $K$ dividing $m$. Let $\hat{K}$ be the maximal abelian extension of $K$ of exponent $m$ and unramified outside $S$. Then $\hat{K}$ is a finite Galois extension of $K$.
	\end{theorem}
	\begin{proof}
		
	\end{proof}
	Now we can easily conclude the proof of \thmref{weakmordellweil}. Since the extensions $K(y)/K$ described in \propref{kummerproperties} are abelian, of exponent $m$ and unramified outside $S$, they all lie inside $\hat{K}$. Since the extension $\hat{K}$ is finite by \thmref{maxabextension}, it has a finite number of subextensions, so there is a finite number of extensions $K(y)$. Since $L$ is, by definition, the compositum of these extensions and they are all finite, $L/K$ must be a finite extension, thus concluding the proof of the weak Mordell-Weil theorem.
	
	\subsection{Height functions}
	
	
	In this section, we will study the so called height functions, which are certain functions constructed for smooth projective varieties that, when defined for an abelian variety, are compatible with its geometry and arithmetics in a way that allows us to measure the complexity of its points. With these functions, we will be able to deduce \thmref{mordellweiltheorem} from \thmref{weakmordellweil}. The use of height functions in the proof of the Mordell-Weil theorem is not entirely necessary in some cases, but since we want to prove the theorem in more generality it is convenient to use them. Furthermore, the theory of height functions is interesting in its own, since it is a key ingredient in many other finiteness theorems from arithmetic geometry.
	
	Rather than following the usual direction of giving definitions and proving theorems immediately after, we will try to explain the reason for each of the definitions and results that we present in a logical order so that each step appears natural to take.
	
	Since we will be dealing with projective varieties, we should first understand the theory of height functions for projective space. Let us fix some notation:
	\begin{itemize}
		\item $K$ will denote a number field, and $K'$ some finite extension of $K$.
		\item $\P$ will denote a prime of $K$, with associated absolute value $|\cdot|_\P$.
		\item An absolute value $|\cdot|$ is called standard if its restriction to $\Q$ is either the absolute value $|\cdot|_\infty$ or one the $p$-adic absolute values $|\cdot|_p$. $M_K$ will denote the set of standard absolute values on $K$, with $M_K^0$ denoting the nonarchimedean ones and $M_K^\infty$ denoting the archimedean ones. 
		\item $K_\P$ will denote the completion of $K$ at $\P$.
		\item If $\P$ lies over the rational prime $p$, the local degree is $n_\P:=[K_\P:\Q_p]$.
		\item The normalized absolute value associated to $\P$ is $||\cdot||_\P=|\cdot|_\P^{n_\P}$.
		\item As in last section, if $S$ is a finite set of nonarchimedean primes of $K$, we denote the ring of $S$-integers of $K$ by $\mathcal{O}_{K,S}$.
	\end{itemize}
	
	With this notation, we have the identity
	\begin{equation}\label{productformula}
		\prod_{\P\in M_K}||x||_\P=1,\text{ for all } x\in K^*,
	\end{equation}
	which is called the \emph{product formula}. If $\P$ is a prime of $K$ and $K'$ is a finite extension of $K$, we also have the identity
	\begin{equation}\label{degreeformula}
		\sum_{\substack{\P\in M_{K'}\\\P'|\P}}[K'_{\P'}:K_\P]=[K':K],
	\end{equation}
	which is called the degree formula.
	
	To gather some intuition, we will begin by the easiest case: the projective space over the rational numbers. We want to define a function on $\mathbb{P}^n(\Q)$ that measures the size of its points. The first problem that we observe is that we should take into consideration the equivalence relation on $\mathbb{P}^n(\Q)$ to define such a function. For now, we can omit this problem if we consider a special set of homogeneous coordinates for our points, using the fact that $\Z=\mathcal{O}_\Q$ is a principal ideal domain. Observe that any point $P\in\mathbb{P}^n(\Q)$ can be written in the form:
	\begin{equation}\label{coordinates}
		P=[x_0:\ldots:x_n],\quad\text{with } x_i\in\Z\quad\text{and } \gcd(x_0,\ldots,x_n)=1,
	\end{equation}
	by clearing denominators and common factors if necessary. Moreover, these $x_i$ are uniquely determined up to sign, so we can define the \emph{height of $P$} to be:
	\begin{equation*}
		H(P)=\max\{|x_0|,\ldots,|x_n|\}.
	\end{equation*}
	This definition has the properties that we want, namely the fact that for any fixed number $C$, the set
	\begin{equation*}
		\{P\in\mathbb{P}^n(\Q): H(P)\le C\}
	\end{equation*}
	is finite, because only a finite number of integers are bounded by $C$. We would like to generalize this definition to arbitrary number fields $K$, but a new problem arises because in general the ring of integers of a number field is not a principal ideal domain. To solve this, we make the following observation: by the assumptions on the $x_i$ above, for any prime $p\in\Z$ there is some $x_j$ with maximal $p$-adic absolute value among the $x_i$. Since the $x_i$ are integers, we have $|x_i|_p\le 1$. If we had the strict inequality $|x_i|<1$ for all $i$, then $p$ would divide the greatest common divisor of the $x_i$, which is impossible. Thus we must have $|x_j|=1$, so we find that the quantity
	\begin{equation*}
		\max\{|x_0|_p,\ldots,|x_n|_p\}
	\end{equation*}
	is equal to one. This motivates the definition of the \emph{height of $P$ (relative to $K$)} as the quantity:
	\begin{equation*}
		H_K(P):=\prod_{\P\in M_K}\max\{||x_0||_\P,\ldots,||x_n||_\P\}.
	\end{equation*}
	Since later we will deal with Abelian varieties and we want to study the behavior of heights with respect to the addition of points, it is convenient to define the \emph{logarithmic height of $P$ (relative to $K$)} as
	\begin{equation*}
		h_K(P):=\log H_K(P).
	\end{equation*}
	Note that now we are not making any assumption on the coordinates of $P\in\mathbb{P}^n(K)$, so we should verify that $H$ is well defined. Moreover, we would like to relate the height $H$ to the relative heights $H_K$. The following proposition clears this.
	\begin{proposition}
		Let $K$ be a number field and let $P\in\mathbb{P}^n(K)$ be a point.
		\begin{enumerate}
			\item The value of $H_K(P)$ is independent of the choice of homogeneous coordinates for $P$.
			\item $H_K(P)$ is greater or equal than 1, for all $P\in\mathbb{P}^n(K)$.
			\item If $K'$ is a finite extension of $K$, then
			\begin{equation*}
				H_{K'}(P)=H_K(P)^{[K':K]}.
			\end{equation*}
		\end{enumerate}
	\end{proposition}
	\begin{proof}
	For the first assertion, let $P=(x_0,\ldots,x_n)$. If we now multiply the coordinates by $c\in K^*$, so that they represent the same point in $\mathbb{P}^n(K)$, we find:
	\begin{align*}
		&\prod_{\P\in M_K}\max\{||cx_0||_\P,\ldots,||cx_n||_\P\}\\
		&=\left(\prod_{\P\in M_K}||c||_\P\right) \left(\prod_{\P\in M_K}\max\{||x_0||_\P,\ldots,||x_n||_\P\}\right)\\
		&=\prod_{\P\in M_K}\max\{||x_0||_\P,\ldots,||x_n||_\P\},
	\end{align*}
	where the last equality follows from the product formula \eqref{productformula}.
	
	For the second assertion, note that we can always find coordinates for $P$ such that some coordinate is equal to 1, and the value of $H_K$ will not change by the previous assertion. Computing $H_K$ in this case we see that it must be greater or equal than 1, since $|1|=1$ for any absolute value $|\cdot|$ of $K$.
	
	For the last assertion, we compute using the definition:
	\begin{align*}
		H_{K'}(P)&=\prod_{\P'\in M_{K'}}\max\{||x_0||_{\P'},\ldots,||x_n||_{\P'}\}\\
		&=\prod_{\P\in M_K}\prod_{\substack{\P'\in M_{K'}\\\P'|\P}}\max\{||x_0||_{\P'},\ldots,||x_n||_{\P'}\}\\
		&=\prod_{\P\in M_K}\prod_{\substack{\P'\in M_{K'}\\\P'|\P}}\max\{|x_0|^{n_{\P'}}_{\P},\ldots,|x_n|^{n_{\P'}}_{\P}\},
	\end{align*}
	where the last equality holds because $|\cdot|_{\P'}$ restricts to $|\cdot|_\P$ on $K$. Since we have $n_{\P'}=[K'_{\P'}:\Q]=[K'_{\P'}:K_\P][K_\P:\Q]=[K'_{\P'}:K_\P]n_\P$, we continue:
	\begin{align*}
		&=\prod_{\P\in M_K}\prod_{\substack{\P'\in M_{K'}\\\P'|\P}}\max\{||x_0||_\P,\ldots,||x_n||_\P\}^{[K'_{\P'}:K_\P]}\\
		&=\prod_{\P\in M_K}\max\{||x_0||_\P,\ldots,||x_n||_\P\}^{[K':K]}\quad(\text{From the degree formula \eqref{degreeformula}})\\
		&=H_K(P)^{[K':K]},
	\end{align*}
	as we wanted.
	\end{proof}
	As a consequence of the last proposition, we observe that $H_\Q(P)$ is independent on the choice of coordinates of $P$, so we can take them as in \eqref{coordinates} and deduce that $H_\Q(P)$ is the same function that we defined at the beginning. Furthermore, we can take $[K':\Q]$-th roots in the identity of part $(c)$ and obtain:
	\begin{equation*}
		H_{K'}^{1/[K':\Q]}(P)=H_{K}^{[K':K]/[K':\Q]}(P)=H_K^{1/[K:\Q]}(P),
	\end{equation*}
	which motivates the following definition: the \emph{absolute (multiplicative) height} on $\mathbb{P}^n$ is the function
	\begin{equation*}
		H:\mathbb{P}^n(\overline{\Q})\rightarrow[1,\infty),\quad H(P)=H_K(P)^{1/[K:\Q]}
	\end{equation*}
	where $K$ is any number field such that $P\in\mathbb{P}^n(K)$. Similarly, we define the \emph{absolute (logarithmic) height} on $\mathbb{P}^n$ as the function
	\begin{equation*}
		h:\mathbb{P}^n(\overline{\Q})\rightarrow[0,\infty),\quad h(P)=\log H(P).
	\end{equation*}
	We also define the height of an element $\alpha\in K$ as the height of the point in the projective line $[\alpha:1]\in\mathbb{P}^1(K)$, whether if we are referring to the absolute or relative height, multiplicative or logarithmic. The following lemma explains the relation between the absolute height function and the absolute Galois group of $\Q$.
	
	\begin{lemma}\label{heightgalois}
		Let $P$ be a point in $\mathbb{P}^n(\overline{\Q})$ and let $\sigma\in G_\Q$. Then $H(P)=H(P^\sigma)$.
	\end{lemma}
	\begin{proof}
		See \cite{Hindry_Silverman_2000}. The idea for the proof is to use the relations between an absolute value $|\cdot|_\P$ and $|\sigma(\cdot)|_{\sigma(\P)}$, for $\sigma\in G_\Q$.
	\end{proof}
	
	Observe that if $P\in\mathbb{P}^n(\Q)$, then $H(P)=H_\Q(P)$, so the absolute height is the proper generalization of our earlier considerations. To further remark the generalization, we prove a finiteness result for the absolute height function. This will be crucial in our later work, since we will be able to reduce in many cases to this situation. Recall that we define the field of definition of a point $P=[x_0,\ldots,x_n]\in\mathbb{P}^n(\overline{\Q})$ as the field
	\begin{equation*}
		\Q(P)=\Q(x_0/x_j,\ldots,x_n/x_j),
	\end{equation*}
	for any $j$ such that $x_j\neq 0$. We then have:
	
	\begin{proposition}[Northcott property]
		Let $C$ and $D$ be any two positive numbers. Then the set
		\begin{equation*}
			\{P\in\mathbb{P}^n(\overline{\Q}): H(P)\le C \text{ and } [\Q(P):\Q]\le D\}
		\end{equation*}
		is finite. In particular, for any fixed number field $K$, the set
		\begin{equation*}
			\{P\in\mathbb{P}^n(K):H_K(x)\le C\}
		\end{equation*}
		is finite.
	\end{proposition}
	\begin{proof}
		We first prove a seemingly weaker assertion: for any $d\ge 1$, the set
		\begin{equation*}
			\{x\in\overline{\Q}: H(x)\le C\text{ and }[\Q(x):\Q]=d\}
		\end{equation*}
		is finite. The strategy is to show that such an element $x$ is the root of a polynomial with bounded coefficients, where the bounds are independent of $x$. Let $x\in\overline{\Q}$ be of degree $d$ and such that $H(x)\le C$. Denote by $k$ the field $\Q(x)$. Let $x_1,\ldots,x_d$ be the Galois conjugates of $x$ over $\Q$, and let $F_x(T)$ be the minimal polynomial of $x$ over $\Q$. By elementary Galois theory, we can write $F_x(T)$ as a combination of the elementary symmetric polynomials:
		\begin{equation*}
			F_x(T)=\prod_{i=1}^d(T-x_j)=\sum_{r=0}^d(-1)^rs_r(x)T^{d-r},
		\end{equation*}
		where
		\begin{equation*}
			s_r(x)=\sum_{1\le i_1\le\ldots\le i_r\le d}x_{i_1}\cdot\ldots\cdot x_{i_r}
		\end{equation*}
		is the $r^\text{th}$ elementary symmetric polynomial. We can then bound, for any absolute value $|\cdot|_\P$:
		\begin{align*}
			|s_r(x)|_\P&=\left|\sum_{1\le i_1\le\ldots\le i_r\le d}x_{i_1}\cdot\ldots\cdot x_{i_r}\right|_\P\\
			&\le c(\P,r,d)\max_{1\le i_1\le\ldots\le i_r\le d}|x_{i_1}\cdot\ldots\cdot x_{i_r}|_\P\\
			&\le c(\P,r,d)\max_{1\le i\le d}|x_i|^r_\P,
		\end{align*}
		where $c(\P,r,d)$ is 1 if $\P$ is nonarchimedean, by the ultrametric inequality; and it is ${d\choose r}\le 2^d$ if $\P$ is archimedean. In either case, $c(\P,r,d)$ can be bounded by a value $c(\P,d)$ which is independent of $r$: 1 if $\P$ is nonarchimedean and $2^d$ if it is archimedean. It is clear then that
		\begin{equation*}
			\max\{|s_0(x)|_\P,\ldots,|s_d(x)|_\P\}\le c(\P,d)\prod_{i=1}^d\max\{|x_i|_\P,1\}^d.
		\end{equation*}
		If we take the product of this inequality over all $\P$ and then take the $[k:\Q]^{\text{th}}$ root, we obtain:
			\begin{equation*}
				H(s_0(x),\ldots,s_d(x))\le 2^d\prod_{i=1}^dH(x_i)^d.
			\end{equation*}
		Now since the height function is invariant by the action of the Galois group on $\mathbb{P}^n(\overline{Q})$ (\lemref{heightgalois}), we see that $H(x_i)=H(x)$ for all $i$, so we deduce:
		\begin{equation*}
			H(s_0(x),\ldots,s_d(x))\le 2^dH(x)^{d^2}\le 2^dC^{d^2}.
		\end{equation*}
		This means that the coefficients of $F_x(T)$, which are in $\Q$, have bounded height. As we have already seen, there are only finitely many possibilities for points in $\mathbb{P}^n(\Q)$ to have bounded height. Therefore, any $x$ with $H(x)\le C$ and degree $d$ will be a root of a polynomial from a finite collection, so there can only be a finite number of such $x$.
		
		Now we finish the proof: let $P$ be a point in $\mathbb{P}^n(\overline{\Q})$ with homogeneous coordinates $P=[x_0:\ldots:x_n]$ such that some coordinate is equal to 1. Assume that $H(P)\le C$ and $[\Q(P):\Q]\le D$. If $|\cdot|_\P$ is an absolute value, then there is some $x_j$ with maximal absolute value among the $x_i$, which will necessarily be greater or equal than one. Since it will also be greater or equal than the other $|x_i|_\P$, we find that for any absolute value $|\cdot|$ and any index $i$,
		\begin{equation*}
			\max\{||x_0||_\P,\ldots,||x_n||_\P\}\ge \max\{{||x_i||_\P,1}\}.
		\end{equation*}
		Multiplying over all the relevant absolute values, we obtain an inequality on the relative height functions $H_K(P)\ge H_K(x_i)$, where $K$ is a field in which $P$ is defined. Taking the $[K:\Q]$-th square root of both sides, we obtain the inequality $H(P)\ge H(x_i)$, for all $i$. In addition, we have that $\Q(P)\supset\Q(x_i)$, so $[\Q(P):\Q]\ge [\Q(x_i):\Q]$. Therefore, each of the coordinates of $P$ lies in one of the following finite collection of sets
		\begin{equation*}
			\{x\in\overline{\Q}: H(x)\le C\text{ and }[\Q(x):\Q]=d\}, \text{ with } 1\le d\le D
		\end{equation*}
		which we have proven to be finite. Therefore, there is only a finite number of possibilities for the coordinates of $P$, and hence there can only be a finite number of such points $P$.
	\end{proof}
	
	We are now ready to pass from the projective space to smooth projective varieties in general. To do that, we will use our knowledge of height functions on $\mathbb{P}^n$. If $V$ is a projective variety defined over $\overline{Q}$, and $\phi:V\rightarrow \mathbb{P}^n$ is a morphism, we can define the \emph{(absolute logarithmic) height on $V$ relative to $\phi$} as:
	\begin{align*}
		h_\phi:V(\overline{K})&\rightarrow [0,\infty]\\
		P&\mapsto h_\phi(P)=h(\phi(P)),
	\end{align*}
	where $h$ is the absolute logarithmic height on $\mathbb{P}^n$.
	
	We will need the following lemma from algebraic geometry.
	\begin{lemma}
		Let $V$ be a projective variety defined over $\overline{Q}$, let $\phi:V\rightarrow\mathbb{P}^n$ and $\psi:V\rightarrow\mathbb{P}^m$ be morphisms, and let $H$ and $H'$ be hyperplanes in $\mathbb{P}^n$ and $\mathbb{P}^m$, respectively. Assume that $\phi^*H$ and $\psi^*H'$ are linearly equivalent. Then
		\begin{equation*}
			h_\phi(P)=h_\psi(P)+O(1), \text{ for all } P\in V(\overline{\Q}),
		\end{equation*}
		where the constant from $O(1)$ depends on $V$, $\phi$, and $\psi$, but not on $P$.
	\end{lemma}
	\begin{proof}
		We refer to \cite{Hindry_Silverman_2000}.
	\end{proof}
	
	With this result, we are going to prove a fundamental theorem, due to André Weil, that assigns height functions to each divisor on a smooth projective variety which is compatible with the different structures of $V,\Div(V)$ and the category of projective varieties.
	
	The smoothness assumption is not actually needed, since one can always use Cartier divisors. However, our main objects of study are abelian varieties, which are smooth, so we will explain the theory only for smooth projective varieties since it will allow the use of the simpler language of Weil divisors.
	
	We have defined heights relative to morphisms to $\mathbb{P}^n$, but in order to assign heights to divisors, we have to explain how to relate divisors on $V$ with morphisms $\phi:V\rightarrow\mathbb{P}^n$. We will recall some terminology from algebraic geometry.
	\begin{definition}
		Let $V$ be a smooth projective variety and $D\in\Div(V)$ a divisor. The complete linear system of $D$, denoted by $|D|$ is the set of all effective divisors linearly equivalent to $D$, together with the structure of projective space given by the parametrization
		\begin{align*}
			\mathbb{P}(L(D))&\rightarrow \{D'\in\Div(V):D\sim D', D'\ge 0\}\\
			f\pmod{K^*}&\mapsto D+\Div(f).
		\end{align*}
				
		A linear system $\mathcal{L}$ is a subset of some complete linear system $|D|$ which is a linear subspace for the projective space structure of $|D|$. In other words, a linear system is the projective space of a vector space $E$, which is a vector subspace of $L(D)$, for some $D$.
		
		The dimension of a linear system is its dimension as a linear projective variety.
	\end{definition}
	To any linear system $\mathcal{L}$, we can associate a rational map $\phi_L$ as follows: $\mathcal{L}$ is parametrized by some projective space $\mathbb{P}(E)\subset \mathbb{P}(L(D))$. Choose a basis $f_0,\ldots, f_n$ of $E\subset L(D)$. Then we define
	\begin{align*}
		\phi_L:V&\dasharrow \mathbb{P}^n\\
		x&\mapsto (f_0(x),\ldots,f_n(x)),
	\end{align*}
	which is a rational map on $X$. A different choice of basis changes the map $\phi_L$ by an automorphism of $\mathbb{P}^n$. Note that in general, $\phi_L$ is not a morphism since it is not defined on the set of common zeros of the $f_i$. The set of base points of $\mathcal{L}$ is the intersection of the supports of all divisors in $\mathcal{L}$. A linear system whose set of base points is empty is called a base point free linear system. We say that a divisor $D$ is base point free if the complete linear system $|D|$ is base point free. The important property of base point free divisors (or linear systems) is that they correspond to morphisms $\phi_L:V\rightarrow\mathbb{P}^n$, so this will be type of divisors that we will use to define heights.
	
	Observe that we constructed the map $\phi_L$ using global sections of $D$, so it is convenient to work with divisors that have sufficiently many global sections.
	\begin{definition}
		A linear system is very ample if the associated morphism $\phi_L$ is an embedding, that is, if $\phi_L$ 
		A divisor $D$ is very ample if the linear system $|D|$ is very ample, and it is ample if some positive multiple of $D$ is very ample.
	\end{definition}
	
	In particular ample and very ample divisors are base point free. It can be seen that ample divisors have, in a sense that can be precisely stated, many global sections. We will use the important fact that any divisor can be written as the difference of two base point free (or even ample) divisors, see \cite{Hartshorne_1977}.
	
	\begin{theorem}[Weil's Height Machine]\label{heightmachine}
		Let $K$ be a number field. For any smooth projective variety $V$ defined over $K$, there is a map
		\begin{equation*}
			h_V:\Div(V)\rightarrow\{\text{functions }V(\overline{K})\rightarrow \mathbb{R}\}
		\end{equation*}
		that has the following properties:
		\begin{itemize}
			\item[(a)](Normalization) Let $H\subset\mathbb{P}^n$ be a hyperplane, and let $h(P)$ be the absolute logarithmic height on $\mathbb{P}^n$. Then
			\begin{equation*}
				h_{\mathbb{P}^n,H}(P)=h(P)+O(1)\text{ for all } P\in\mathbb{P}^n(\overline{K}).
			\end{equation*}
			\item[(b)](Functoriality) Let $\phi:V\rightarrow W$ be a morphism and let $D\in\Div(W)$. Then
			\begin{equation*}
				h_{V,\phi^*D}(P)=h_{W,D}(\phi(P))+O(1)\text{ for all }P\in V(\overline{K}).
			\end{equation*}
			\item[(c)](Additivity) Let $D,E\in\Div(V)$. Then
			\begin{equation*}
				h_{V,D+E}(P)=h_{V,D}(P)+h_{V,E}(P)+O(1)\text{ for all }P\in V(\overline{K}).
			\end{equation*}
			\item[(d)](Linear Equivalence) Let $D,E\Div(V)$ be linearly equivalent divisors. Then
			\begin{equation*}
				h_{V,D}(P)=h_{V,E}(P)+O(1)\text{ for all }P\in V(\overline{K}).
			\end{equation*}
			\item[(e)](Positivity) Let $D\Div(V)$ be an effective divisor, and let $B$ be the base locus of the linear system $|D|$. Then
			\begin{equation*}
				h_{V,D}\ge O(1)\text{ for all }P\in (V\setminus B)(\overline{K}).
			\end{equation*}
			\item[(f)](Finiteness) Let $D\Div(V)$ be ample. Then for every finite extension $K'/K$ and every constant $B$, the set
			\begin{equation*}
				\{P\in V(K'):h_{V,D}(P)\le B\}
			\end{equation*}
			is finite.
			\item[(g)](Uniqueness) The height functions $h_{V,D}$ are determined, up to $O(1)$, by properties (a), (b) just for embeddings $\phi:V\hookrightarrow\mathbb{P}^n$, and (c).
		\end{itemize}
	\end{theorem}
	\begin{proof}
		
	\end{proof}
	In addition to the construction of height functions, the strength of Weil's height machine is that it allows to translate geometric relations of divisors to arithmetic relations on height functions. For example, if $A$ is an abelian variety over a number field $K$ and $m$ is an integer, Mumford's formula explains the behavior of divisors with respect to the map $[m]:A\rightarrow A$:
	\begin{equation*}
		[m]^*D\sim\frac{m^2+m}{2}D+\frac{m^2-m}{2}[-1]^*D.
	\end{equation*}
	(see appendix). Using the height machine, we can deduce important facts about the height functions of $A$.
	
	\begin{corollary}\label{mumfordheight}
		Let $A$ be an abelian variety defined over a number field $K$, and let $D\Div(A)$ be a divisor.
		\begin{itemize}
			\item[(a)] Let $m$ be an integer. Then for all $P\in A(\overline{K})$, we have
			\begin{equation*}
				h_{A,D}([m]P)=\frac{m^2+m}{2}h_{A,D}(P)+\frac{m^2-m}{2}h_{A,D}(-P)+O(1).
			\end{equation*}
			In particular, if $D$ has a symmetric divisor class, we have:
			\begin{equation*}
				h_{A,D}([m]P)=m^2h_{A,D}(P)+O(1).
			\end{equation*}
			\item[(b)] If $D$ has a symmetric divisor class, then for all $P,Q\in A(\overline{K})$, we have:
			\begin{equation*}
				h_{A,D}(P+Q)+h_{A,D}(P-Q)=2h_{A,D}(P)+2h_{A,D}(Q)+O(1).
			\end{equation*}
		\end{itemize}
	\end{corollary}
	\begin{proof}
		(a) From Mumford's formula
		\begin{equation*}
			[m]^*D\sim\frac{m^2+m}{2}D+\frac{m^2-m}{2}[-1]^*D,
		\end{equation*}
		we use the properties of Weil's height machine to compute:
		\begin{align*}
			h_{A,D}([m]P)&=h_{A,[m]^*D}(P)+O(1)\\
			&=h_{A,(1/2)(m^2+m)D+(1/2)(m^2-m)[-1]^*D}(P)+O(1)\\
			&=\frac{m^2+m}{2}h_{A,D}(P)+\frac{m^2-m}{2}h_{A,[-1]^*D}(P)+O(1)\\
			&=\frac{m^2+m}{2}h_{A,D}(P)+\frac{m^2-m}{2}h_{A,D}(-P)+O(1),
		\end{align*}
		as we wanted. If $D$ has a symmetric divisor class, Mumford's formula reads
		\begin{equation*}
			[m]^*D\sim m^2D,
		\end{equation*}
		so we have 
		\begin{align*}
			h_{A,D}([m]P)&=h_{A,[m]^*D}(P)+O(1)\\
			&=h_{A,m^2D}(P)+O(1)\\
			&=m^2h_{A,D}(P)+O(1).
		\end{align*}
		(b) If $D$ has a symmetric divisor class, consider the maps $s_{12},d_{12},p_1,p_2:A\times A\rightarrow A$ defined as
		\begin{equation*}
			s_{12}(P,Q)=P+Q,\quad d_{12}(P,Q)=P-Q,\quad p_1(P,Q)=P,\quad\text{and }p_2(P,Q)=Q.
		\end{equation*}
		Then we have the relation
		\begin{equation*}
			s_{12}^*D+d_{12}^*D=2p_1^*D+2p_2^*D
		\end{equation*}
		(see appendix). Using the properties of Weil's height machine, we find:
		\begin{align*}
			h_{A,D}(P+Q)+h_{A,D}(P-Q)&=h_{A,D}(s_{12}(P,Q))+h_{A,D}(d_{12}(P,Q))\\
			&=h_{A,s_{12}^*D}(P,Q)+h_{A,d_{12}^*D}(P,Q)+O(1)\\
			&=h_{A,s_{12}^*D+d_{12}^*D}(P,Q)+O(1)\\
			&=h_{A,2p_1^*D+2p_2^*D}(P,Q)+O(1)\\
			&=h_{A,2p_1^*D}(P,Q)+h_{A,2p_2^*D}(P,Q)+O(1)\\
			&=2h_{A,D}(p_1(P,Q))+2h_{A,D}(p_2(P,Q))+O(1)\\
			&=2h_{A,D}(P)+2h_{A,D}(Q)+O(1)
		\end{align*}
		which is what we wanted.
	\end{proof}
	
	The height machine gives us relations between the height functions up to bounded functions. In order to prove deeper results, we would like to get rid of the $O(1)$. André Néron and John Tate studied this problem, and discovered that under certain assumptions on $D$ being an eigenvalue of the pullback of a morphism $\phi:V\rightarrow V$, it is possible to find (and compute) a height function on the $O(1)$ class of $h_{V,D}$ with useful properties.
	
	\begin{theorem}[Néron,Tate]\label{canonicalheight}
		Let $V$ be a smooth variety defined over a number field $K$, let $D\in\Div(A)$ be a divisor, and let $\phi:V\rightarrow V$ be a morphism. Suppose that there exists a number $\alpha>1$ such that
		\begin{equation*}
			\phi^*D\sim\alpha D.
		\end{equation*}
		Then there is a unique function, called the canonical height on $V$ relative to $\phi$ and $D$,
		\begin{equation*}
			\hat{h}_{V,\phi,D}:V(\overline{K})\rightarrow\mathbb{R},
		\end{equation*}
		such that:
		\begin{itemize}
			\item[(a)]$\hat{h}_{V,\phi,D}(P)=h_{V,D}(P)+O(1)$, for all $P\in V(\overline{K})$.
			\item[(b)]$\hat{h}_{V,\phi,D}(\phi(P))=\alpha\hat{h}_{V,\phi,D}(P)$, for all $P\in V(\overline{K})$.
		\end{itemize}
		The canonical height depends only on the linear equivalence class of $D$, and it can be computed as the limit
		\begin{equation*}
			\hat{h}_{V,\phi,D}(P)=\lim\limits_{n\to\infty}\frac{1}{\alpha^n}h_{V,D}(\phi^n(P)),
		\end{equation*}
		where $\phi^n$ is the $n^\text{th}$ iterate of $\phi$.
	\end{theorem}
	\begin{proof}
		Applying the Height machine to the relation $\phi^*D\sim\alpha D$, we obtain that:
		\begin{equation*}
			h_{V,D}(\phi(P))=\alpha h_{V,D}(P)+O(1), \text{ for all } P\in V(\overline{K}).
		\end{equation*}
		As always, the constant $C$ bounding the difference of the functions above depends on $V,\phi,D$, and the choice of $h_{V,D}$, but not on $P$. Now we will show that the sequence $\frac{1}{\alpha^n}h_{V,D}(\phi^n(P))$ from the statement of the theorem is a Cauchy sequence, so it converges. Let $n\ge m$, then:
		
		\begin{align*}
			|\alpha^{-n}h_{V,D}(\phi^n(&P))-\alpha^{-m}h_{V,D}(\phi^m(P))|\\
			&=\left|\sum_{i=m+1}^n\alpha^{-i}(h_{V,D}(\phi^i(P))-\alpha h_{V,D}(\phi^{i-1}(P)))\right|\\
			&\le\sum_{i=m+1}^n\alpha^{-i}\left|h_{V,D}(\phi^i(P))-\alpha h_{V,D}(\phi^{i-1}(P))\right|\\
			&\le\sum_{i=m+1}^n\alpha^{-i}C\\
			&=\left(\frac{\alpha^{-m}-\alpha^{-n}}{\alpha-1}\right)C.
		\end{align*}
		As $n>m$ goes to infinity, the last expression goes to zero so our sequence is Cauchy. We can then define $\hat{h}_{V,\phi,D}$ as
		\begin{equation*}
			\hat{h}_{V,\phi,D}(P)=\lim\limits_{n\to\infty}\frac{1}{\alpha^n}h_{V,D}(\phi^n(P)),
		\end{equation*}
		as we claimed. Now we just verify the properties. For the first one, take $m=0$ and let $n\to\infty$ in the above inequality. This gives
		\begin{equation*}
			|\hat{h}_{V,\phi,D}(P)-h_{V,D}(P)|\le\frac{1}{\alpha-1}C,
		\end{equation*}
		which proves (a). To prove (b), we just use the definition:
		\begin{align*}
			\hat{h}_{V,\phi,D}(\phi(P))&=\lim\limits_{n\to\infty}\frac{1}{\alpha^n}h_{V,D}(\phi^n(\phi(P)))\\
			&=\lim\limits_{n\to\infty}\frac{\alpha}{\alpha^{n+1}}h_{V,D}(\phi^{n+1}(P))\\
			&=\alpha\hat{h}_{V,\phi,D}(P).
		\end{align*}
		It only remains to show that properties (a) and (b) determine $\hat{h}_{V,\phi,D}$. Suppose that there are two functions $\hat{h}$, $\hat{h}'$ satisfying (a) and (b). Then, by property (a) we have:
		\begin{align*}
			\hat{h}(P)-\hat{h}'(P)=h_{V,D}(P)-h_{V,D}(P)+O(1)=O(1),\text{ for all }P\in V(\overline{K}).
		\end{align*}
		Therefore the quantity $|\hat{h}-\hat{h}'|$ is bounded by a certain constant $C_1$. Since $(\hat{h}-\hat{h}')\circ\phi=\alpha(\hat{h}-\hat{h}')$ by property (b), we have that $(\hat{h}-\hat{h}')\circ\phi^n=\alpha^n(\hat{h}-\hat{h}')$. Therefore
		\begin{equation*}
			|(\hat{h}-\hat{h}')(P)|=\frac{|(\hat{h}-\hat{h}')\circ\phi^n|}{\alpha^n}\le \frac{C_1}{\alpha^n},
		\end{equation*}
		and the last quantity goes to zero as $n$ goes to infinity, so $\hat{h}-\hat{h}'=0$, as we wanted.
	\end{proof}
	
	In order to use \ref{canonicalheight} to produce canonical height functions, we have to find a morphism $\phi:V\rightarrow V$ and a divisor $D$ that is an eigenvalue for $\phi^*$, up to linear equivalence. Recall from \corref{mumfordheight} that in an abelian variety $A$, symmetric divisors are, up to linear equivalence, eigenvalues for all the maps $[m]$: $[m]^*D\sim m^2D$. This will make it possible to find a canonical height on $A$ relative to $D$ and to all the maps $[m]$ simultaneously.
	
	\begin{theorem}\label{nerontate}
		Let $A$ be an abelian variety defined over some number field $K$, and let $D\in\Div(A)$ be a divisor in a symmetric divisor class. There is a height function
		\begin{equation*}
			\hat{h}_{A,D}:A(\overline{K})\rightarrow\mathbb{R},
		\end{equation*}
		called the canonical height on $A$ relative to $D$, with the properties:
		\begin{itemize}
			\item[(a)]$\hat{h}_{A,D}(P)=h_{A,D}(P)+O(1)$, for all $P\in A(\overline{K})$.
			\item[(b)]For all integers $m$,
			\begin{equation*}
				\hat{h}_{A,D}([m]P)=m^2\hat{h}_{A,D}(P),\text{ for all }P\in A(\overline{K}).
			\end{equation*}
			\item[(c)](Parallelogram law)
			\begin{equation*}
				\hat{h}_{A,D}(P+Q)+\hat{h}_{A,D}(P-Q)=2\hat{h}_{A,D}(P)+2\hat{h}_{A,D}(Q),\text{ for all }P,Q\in A(\overline{K}).
			\end{equation*}
			\item[(d)](Uniqueness) The canonical height $\hat{h}_{A,D}$ depends only on the divisor class of the divisor $D$. It is uniquely determined by $(a)$ and $(b)$ for any single integer $m\ge 2$.
		\end{itemize}
	\end{theorem}
	\begin{proof}
		Recall that for a symmetric divisor $D\Div(A)$ and an integer $r$, we have $[r]^*D\sim r^2D$. Therefore, we may use \thmref{canonicalheight} with any $r\ge 2$, so we will use $r=2$. We thus obtain the canonical height on $A$ relative to $[2]$ and $D$, which we may compute as the limit
		\begin{equation*}
			\hat{h}_{A,[2],D}(P)=\lim\limits_{n\to\infty}\frac{1}{4^n}h_{A,D}([2^n]P).
		\end{equation*}
		The first assertion is just property (a) of the canonical height. Note that the second assertion for $m=2$ is just property (b) of the canonical height. Since canonical heights are determined by these two properties, when we prove the second assertion of this theorem we will actually have shown that $\hat{h}_{A,[2],D}$ is canonical relative to every map $[m]$ (and thus we will omit the $[2]$ in the notation).
		
		Let $m\ge2$ be any integer. Then we may compute, using the definition:
		\begin{align*}
			\hat{h}_{A,[2],D}([m]P)&=\lim\limits_{n\to\infty}\frac{1}{4^n}h_{A,D}([2^n][m]P)\\
			&=\lim\limits_{n\to\infty}\frac{1}{4^n}h_{A,D}([m][2^n]P)\\
			&=\lim\limits_{n\to\infty}\frac{1}{4^n}\left(h_{A,m^2D}([2^n]P)+O(1)\right)\\
			&=\lim\limits_{n\to\infty}\frac{1}{4^n}\left(m^2h_{A,D}([2^n]P)+O(1)\right)\\
			&=m^2\lim\limits_{n\to\infty}\frac{1}{4^n}h_{A,D}([2^n]P)=m^2\hat{h}_{A,[2],D}(P),
		\end{align*}
		which proves the second assertion. As we explained, we may now update the notation to $\hat{h}_{A,D}:=\hat{h}_{A,[2],D}$. For the third assertion, we use the relation:
		\begin{align*}
			h_{A,D}(P+Q)+h_{A,D}(P-Q)=2h_{A,D}(P)+2h_{A,D}(Q)+O(1)
		\end{align*}
		from \corref{mumfordheight}. Using the definition of $\hat{h}_{A,D}$, we compute:
		\begin{align*}
			\hat{h}_{A,D}(P+Q)+\hat{h}_{A,D}(P-Q)&=\lim\limits_{n\to\infty}\frac{1}{4^n}\left(h_{A,D}([2^n]P+[2^n]Q)+h_{A,D}([2^n]P-[2^n]Q)\right)\\
			&=\lim\limits_{n\to\infty}\frac{1}{4^n}\left(2h_{A,D}([2^n]P)+2h_{A,D}([2^n]Q)+O(1)\right)\\
			&=2\hat{h}_{A,D}(P)+2\hat{h}_{A,D}(Q).
		\end{align*}
		Finally, uniqueness follows from \thmref{canonicalheight} together with the fact that $\hat{h}_{A,D}$ is canonical with respect to every map $[m]$.
	\end{proof}

	The following general lemma will allow us to deduce even more properties for the canonical height associated to a symmetric ample divisor, which will be crucial in the final step of the proof of the Mordell-Weil theorem.
	\begin{lemma}\label{parallelogram}
		Let $A$ be an abelian group, and let $h:A\rightarrow\mathbb{R}$ be a function satisfying the parallelogram law:
		\begin{equation*}
			h(P+Q)+h(P-Q)=2h(P)+2h(Q),\quad\text{ for all } P,Q\in A.
		\end{equation*}
		Then $h$ is a quadratic form on $A$, by which we mean that $h$ is an even function such that the pairing
		\begin{align*}
			\langle\cdot,\cdot\rangle &:A\times A\rightarrow\mathbb{R},\\
			\langle P,Q\rangle&=\frac{h(P+Q)-h(P)-h(Q)}{2}
		\end{align*}
		is bilinear. Furthermore, assume that $h$ is non negative. Defining $|x|:=\sqrt{h(x)}$ for all $x\in A$, we have that $|\cdot|$ is a seminorm.
	\end{lemma}
	\begin{proof}
		Putting $P=Q=0$ in the formula for the parallelogram law we obtain that $2h(0)=4h(0)$, so $h(0)=0$. Now putting $P=0$ in the formula, we obtain $h(Q)+h(-Q)=2h(Q)$, so $h(Q)=h(-Q)$. We now have to show that the associated pairing is bilinear. Since it is symmetric, it suffices to show that it is bilinear in the first component. The equality:
		\begin{equation*}
			\langle P+Q,R\rangle=\langle P,R\rangle+\langle Q,R\rangle
		\end{equation*}
		for all $P,Q,R\in A$ translates into the identity:
		\begin{equation*}
			h(P+Q+R)-h(P+Q)-h(P+R)-h(Q+R)+h(P)+h(Q)+h(R)=0,
		\end{equation*}
		so we have to prove this. Applying the parallelogram law to the pairs $(P+R,Q)$, $(P,R-Q)$, $(P+Q,R)$ and $(R,Q)$ we obtain the following four identities:
		\begin{align*}
			h(P+R+Q)+h(P+R-Q)-2h(P+R)-2h(Q)=&0,\\
			h(P+R-Q)+h(P+Q-R)-2h(P)-2h(R-Q)=&0,\\
			h(P+Q+R)+h(P+Q-R)-2h(P+Q)-2h(R)=&0,\\
			h(R+Q)+h(R-Q)-2h(R)-2h(Q)=&0.
		\end{align*}
		Now adding the first and third equations and substracting the second and fourth ones yields the result.
		
		Assume now that $h$ is non-negative. We first show absolute homogeneity of $|\cdot|=\sqrt{h(\cdot)}$, that is, $|mx|=|m|\cdot|x|$ for any integer $m$ and $x\in A$ (here $|m|$ denotes the usual absolute value). For $m=1$, the condition is trivial. Now we can do induction on $m$: assume that $|mx|=|m|\cdot|x|$ for all $x\in A$. Then we have, for all $x\in A$:
		\begin{equation*}
			\langle (m+1)x,-x\rangle=\frac{h(mx)-h((m+1)x)-h(x)}{2}=\frac{(m^2-1)h(x)-h((m+1)x)}{2}.
		\end{equation*}
		On the other hand, linearity of $\langle\cdot,\cdot\rangle$ implies that
		\begin{equation*}
			\langle (m+1)x,-x\rangle=(m+1)\langle x,-x\rangle=\frac{m+1}{2}[-2h(x)].
		\end{equation*}
		Equating these two terms and clearing denominators yields:
		\begin{equation*}
			h((m+1)x)=(m^2+2m+1)h(x)=(m+1)^2h(x),
		\end{equation*}
		so taking square roots yields $|(m+1)x|=|m+1|\cdot|x|$. Since $h$ is even, absolute homogeneity for negative values of $m$ automatically follows from this. The fact that $|\cdot|$ is non negative follows trivially from the assumption that $h$ is non negative. To prove the triangle inequality, note that the assertion:
		\begin{equation*}
			|x+y|\le |x|+|y|,\quad\text{for all } x,y\in A
		\end{equation*}
		is equivalent, by squaring both sides, to the assertion:
		\begin{equation*}
			h(x+y)\le h(x)+h(y)+2|x||y|,
		\end{equation*}
		which is equivalent to
		\begin{equation*}
			\langle x,y\rangle\le |x||y|.
		\end{equation*}
		And the latter is true by the Cauchy-Schwarz inequality for semi-definite Hermitian forms.
	\end{proof}
	
	There is an analogous theory for canonical height functions relative to antisymmetric divisor classes, which we will not need for our purposes. We state it here for completeness.
	\begin{theorem}
		Let $A$ be an abelian variety defined over a number field $K$, and let $D\in\Div(A)$ be a divisor in an antisymmetric divisor class. There is a canonical height function:
		\begin{equation*}
			\hat{h}_{A,D}:A(\overline{K})\rightarrow\mathbb{R}
		\end{equation*}
		such that:
		\begin{itemize}
			\item[(a)] $\hat{h}_{A,D}(P)=h_{A,D}(P)+O(1)$, for all $P\in A(\overline{K})$.
			\item[(b)] $\hat{h}_{A,D}(P+Q)=\hat{h}_{A,D}(P)+\hat{h}_{A,D}(Q)$, for all $P,Q\in A(\overline{K})$.
		\end{itemize}
		And properties $(a),(b)$ uniquely determine $\hat{h}_{A,D}$.
	\end{theorem}
	Note that for antisymmetric classes, instead of a quadratic form we obtain a homomorphism of groups. Since the double of any divisor can be written as the sum of a symmetric divisor and an antisymmetric divisor, a similar theorem can be obtained for arbitrary divisors. The resulting canonical function is then the sum of a quadratic form and a homomorphism of groups.
		
	Finally, we can conclude the proof of \thmref{mordellweiltheorem} using the weak Mordell-Weil theorem and the theory of heights that we just developed:
	\begin{proof}(of \thmref{mordellweiltheorem})
		
		Since $A$ is a projective variety, we can find an ample divisor $\tilde{D}$ on $A$. Since the sum of ample divisors is ample and the pullback of an ample divisor by a finite morphism is ample, we can consider the divisor
		\begin{equation*}
			D:=\tilde{D}+[-1]^*\tilde{D},
		\end{equation*}
		which is again ample. Furthermore, we see that:
		\begin{equation*}
			[-1]^*D\sim[-1]^*\tilde{D}+[-1\cdot -1]^*\tilde{D}=[-1]^*\tilde{D}+\tilde{D}=D,
		\end{equation*}
		so it is symmetric. Let $\hat{h}$ be the associated Néron-Tate height function on $A(K)$. By \thmref{nerontate} and \lemref{parallelogram}, this function is a nonnegative quadratic form on $A(K)$ such that for every constant $C>0$, the set
		\begin{equation*}
			\{x\in A(K):\hat{h}(x)\le C\}
		\end{equation*}
		is finite. By \thmref{weakmordellweil}, the quotient group $A(K)/mA(K)$ is finite, so we choose a complete set of coset representatives $y_1,\ldots,y_s$. Instead of using directly the height function $\hat{h}$, we will use the function $|x|:=\sqrt{\hat{h}(x)}$ which, thanks to \lemref{parallelogram}, will turn out to be more useful. Let $C_0:=\max_i\hat{h}(y_i)$ and $c_0:=\max_i |y_i|$. Then we claim that the finite set
		\begin{equation*}
			S:=\{x\in A(K):\hat{h}(x)\le C_0\}=\{x\in A(K):|x|\le c_0\}
		\end{equation*}
		generates the group $A(K)$. Indeed, let $x_0\in A(K)$. If $|x_0|\le c_0$, then there is nothing to prove. Assume then that $|x_0|>c_0$. Let $y_i$ be the coset representative for $x_0$ in $A(K)/mA(K)$. Then, $x_0=y_i+mx_1$, for some $x_1\in A(K)$. Using the properties of $\hat{h}$ and \lemref{parallelogram}, we can compute:
		\begin{align*}
			m|x_1|&=\sqrt{m^2\hat{h}(x_1)}=\sqrt{\hat{h}(mx_1)}=|mx_1|\\
			&=|x_0-y_i|\\
			&\le |x_0|+|y_i| \quad(\text{By the triangle inequality})\\
			&< 2|x_0|\quad(\text{Since } |y_i|\le c_0< |x_0|).
		\end{align*}
		Since we are assuming that $m\ge 2$, we find that $|x_1|<|x_0|$. If $x_1$ is in $S$, then we have that $x_0=y_i+mx_1$ is in the group generated by $S$. Otherwise, let $y_j$ be the coset representative for $x_1$ in $A(K)/mA(K)$, so $x_1=y_j+mx_2$, for some $x_2\in A(K)$. By the exact same argument as before, we show that $|x_2|<|x_1|$. Repeating this argument, we obtain a sequence of elements $x_k\in A(K)$ satisfying that $|x_{k+1}|<|x_k|$, and such that $x_k$ is in the group generated by $S$ and $x_{k+1}$. Since there is a finite number of elements with bounded height $\hat{h}$, there exists some $l$ such that $x_l$ is in $S$. Therefore, we find that $x_0$ is in the group generated by $S$. Since $x_0$ was arbitrary, this concludes the proof of the Mordell-Weil theorem.
	\end{proof}
	\newpage
	
	
	
	
	
	
	
	
	
	
	
	
	
	
	
	
	
	\section{Computing the Selmer group}
	
	We have proved in the previous chapter the Mordell Weil theorem for abelian varieties over number fields, so in particular we will turn our attention first to Jacobians of hyperelliptic curves over number fields, and then to Jacobians of hyperelliptic curves over the rational numbers. Given a hyperelliptic curve $C$, we know then that there exists a finite set of points $P_1,\ldots,P_r\in\Jac(C)$ such that every point in $\Jac(C)$ is, up to torsion, a linear combinations of the points $P_i$. However, the proof of this theorem does not give any bound on the rank $r$, nor does it yield any method to find these points.
	
	In this section, we are going to explain a general construction from group cohomology called the Selmer group that allows, in some cases, to compute explicitly (or bound) the rank of a Jacobian variety. Then, we will describe a method for the computation of this group, given by Schaefer in \cite{Schaefer_1995}. This will make it possible for us to bound the Mordell Weil rank of a large amount of (Jacobians of) hyperelliptic curves, in the next chapter.
	
	\subsection{The Selmer group of a hyperelliptic curve}
	
	Let $J$ be a Jacobian variety over a number field $K$. When we proved \thmref{mordellweiltheorem}, we showed that using the theory of height functions it suffices to prove that the weak Mordell Weil group
	\begin{equation*}
		J(K)/mJ(K)
	\end{equation*}
	is finite, for some $m>1$. Moreover, the cardinality of this group gives the rank of $J(K)$: suppose that
	\begin{equation*}
		J(K)\cong\Z^r\oplus J(K)_{\text{tors}},
	\end{equation*}
	where $J(K)_{\text{tors}}$ is the torsion subgroup, which is finite. Then we have
	\begin{equation*}
		mJ(K)=(m\Z)^r\oplus J(K)[m'],
	\end{equation*}
	where $J(K)[m']$ is the subgroup of prime to $m$ torsion. Finally, we obtain that
	\begin{equation*}
		J(K)/mJ(K)\cong (\Zmod{m})^r\oplus J(K)[m].
	\end{equation*}
	Since the torsion is in general easy to compute, being able to determine (the cardinal of) the quotient group $J(K)/mJ(K)$ would give us the exact value of the rank. Unfortunately, it is difficult to work in this quotient group without knowing the group structure of $J(K)$.
	
	Instead, the standard approach is to use the Galois cohomology of the curve to embed the weak Mordell Weil group into a cohomology group, which is finite and effectively computable, and thus provides explicit bounds on the Mordell Weil rank.
	
	The construction is the following: let $J$ be a Jacobian variety over a number field $K$, and let $m>1$ be an integer. Fix once and for all an algebraic closure $\overline{K}$ of $K$. Consider the isogeny:
	\begin{equation*}
		[m]:J\rightarrow J
	\end{equation*}
	given by multiplication by $m$. We denote by $J[m]$ the kernel of this map. Then we have the following short exact sequence:
	\begin{equation}\label{kummersequence}
		\begin{tikzcd}
			0\arrow{r}
			&J[m]\arrow{r}
			&J(\overline{K})\arrow{r}{[m]}
			&J(\overline{K})\arrow{r}
			&0.
		\end{tikzcd}
	\end{equation}
	We denote by $G_K$ the absolute Galois group of $K$, that is, the group $\Gal(\overline{K}/K)$. Note that the above exact sequence of abelian groups is actually an exact sequence of $G_K$ modules, where the action of $G_K$ is the usual one on the coordinates of the points in $J$. Note that the elements that are fixed by $G_K$ are precisely:
	\begin{align*}
		J[m]^{G_K}&=J(K)[m],\\
		J(\overline{K})^{G_K}&=J(K),
	\end{align*}
	so by the usual properties of Galois cohomology we obtain a long exact sequence of the form:
	\begin{equation*}
		\begin{tikzcd}
			0\arrow{r}
			&J(K)[m]\arrow{r}
			&J(K)\arrow{r}{[m]}
			&J(K)\arrow{r}{\delta}
			&\phantom{ }
			\\ \arrow{r}{\delta}
			&H^1(G_K,J[m])\arrow{r}
			&H^1(G_K,J(\overline{K}))\arrow{r}
			&H^1(G_K,J(\overline{K}))\arrow{r}
			& \ldots,
		\end{tikzcd}
	\end{equation*}
	where $\delta$ is the usual connecting homomorphism defined as follows: if $x\in J(K)$ then there exists some $y\in J(\overline{K})$ such that $[m]y=x$. Observe that, since the map $[m]$ is a homomorphism of $G_K$ modules, we have, for each $\sigma\in G_K$:
	\begin{equation*}
		[m](y^\sigma-y)=([m]y)^\sigma -[m]y=x^\sigma -x=x-x=0,
	\end{equation*}
because $x$ is $G_K$ invariant. Therefore, $y^\sigma-y$ is a point of $m$-torsion. Moreover, note that for any $\sigma,\tau\in G_K$ we have:
\begin{equation*}
	y^{\sigma\tau}-y=y^{\sigma\tau}-y^\tau+y^\tau -y=(y^\sigma-y)^\tau + (y^\tau-y),
\end{equation*}
so that the assignment $\sigma\mapsto y^\sigma-y$ has the property of a cocycle. It is a standard exercise in homological algebra to check that this assignment is well defined and does not depend on the choice of $y$, up to coboundaries, so it defines an element in $H^1(G_K, J[m])$, which we define to be $\delta(x)$.

By exactness of the long exact sequence, the kernel of $\delta$ is the image of $[m]:J(K)\rightarrow J(K)$, which is denoted by $mJ(K)$. Therefore, the map $\delta$ induces an injective map 
\begin{equation*}
	\overline{\delta}:J(K)/mJ(K)\hookrightarrow H^1(G_K,J[m]),
\end{equation*}
which we also call $\delta$, by abuse of notation. In addition, the homomorphism $H^1(G_K,J(\overline{K}))\rightarrow H^1(G_K,J(\overline{K}))$ induced by the map $[m]$ can be identified with the multiplication-by-$m$ morphism of the abelian group $H^1(G_K,J(\overline{K}))$, so the kernel of this map, which is the image of the map $H^1(G_K,J[m])\rightarrow H^1(G_K,J(\overline{K}))$, is just $H^1(G_K,J(\overline{K}))[m]$, the $m$ torsion of this group. The above discussion can be summarized by stating that we have a short exact sequence of $G_K$ modules of the form:
\begin{equation}\label{funses}
	\begin{tikzcd}
		0\arrow{r}
		&J(K)/mJ(K)\arrow{r}{\delta}
		&H^1(G_K,J[m])\arrow{r}
		&H^1(G_K,J(\overline{K}))[m]\arrow{r}
		&0.
	\end{tikzcd}
\end{equation}

Since the homomorphism $H^1(G_K,J[m])\rightarrow H^1(G_K,J(\overline{K}))[m]$ is the one induced by the inclusion $J[m]\subset J(\overline{K})$, we will refer to it when needed as $\iota$. Now let $\P$ be a prime of $K$. Denote by $K_\P$, $\overline{K_\P}$ and $G_\P$ the completion of $K$ at $\P$, an algebraic closure of this field and the group $\Gal(\overline{K_\P}/K_\P)$, respectively. Restricting a morphism $\sigma\in G_\P$ to $\overline{K}$ yields an element in $G_K$, and since $\sigma$ fixes $K_\P$, its restriction fixes the prime ideal $\P$. We therefore find that the restriction sends $G_\P$ to the decomposition group $D_\P:=\{\sigma\in G_K: \sigma(\P)=\P\}\subset G_K$. It is easy to see that this is in fact an isomorphism (see \cite{Janusz_1973}), so we can regard the group $G_\P$ as a subgroup of $G_K$. If $M$ is any $G_K$ module, restricting the domain of a 1-cochain $\xi:G_K\rightarrow M$ to $G_\P$ yields the restriction map:
\begin{align*}
	res_\P:H^1(G_K,M)&\rightarrow H^1(G_\P,M)\\
	[\xi]&\mapsto [\xi|_{G_\P}]
\end{align*}
Now we consider all the primes of $K$ simultaneously: for each prime $\P$ of $K$, we can form an analogous short exact sequence to \eqref{funses}, considering $K_\P$ instead of $K$. Taking the product over all primes of $K$ of these exact sequences and using the restriction maps from above, we obtain the following commutative diagram of abelian groups:
\small
	\begin{equation}\label{seldiag}
		\begin{tikzcd}
			0\arrow[r]
			&J(K)/mJ(K)\arrow{r}{\delta}\arrow{d}
			&H^1(G_K, J[m])\arrow{r}\arrow{d}{\prod\text{res}_\P}
			&H^1(G_K,J(\overline{K}))[m]\arrow{r}\arrow{d}{\prod\text{res}_\P}
			&0
			\\
			0\arrow[r]
			&\displaystyle \prod_\P J(K_\P)/mJ(K_\P)\arrow{r}{\delta_\P}
			&\displaystyle \prod_\P H^1(G_\P, J[m])\arrow{r}
			&\displaystyle \prod_\P H^1(G_\P,J(\overline{K}_\P))[m]\arrow{r}
			&0
		\end{tikzcd}
	\end{equation}
\normalsize
	As before, we refer to the homomorphisms $H^1(G_\P,J[m])\rightarrow H^1(G_\P,J)[m]$ as $\iota_\P$ when needed. Now we make the following important observation: let $x\in J(K)/mJ(K)$. Then by commutativity of the above diagram, we have that $\iota_\P\circ\text{res}_\P(\delta(x))$ is equal to $\text{res}_\P\circ \iota(\delta(x))$. Moreover, since the first row is exact, we have that $\iota\circ\delta=0$, so $\text{res}_\P(\delta(x))\in \Im(\delta_\P)$ by exactness of the second row. This motivates the following
	\begin{definition}
		Let $J$ be a Jacobian variety over a number field $K$. The $m$-Selmer group of $J(K)$ is the subgroup of $H^1(G_K,J[m])$ defined as:
		\begin{equation*}
			S^m(J/K)=\{\phi\in H^1(G_K,J[m]):\text{res}_\P(\phi)\in \Im(\delta_\P), \ \forall \P\in\text{Spec}(\mathcal{O}_K)\}
		\end{equation*}
		By exactness of the second row in \eqref{seldiag}, the $m$-Selmer group admits the following equivalent description:
		\begin{equation*}
			S^m(J/K)=\{\phi\in H^1(G_K,J[m]):\text{res}_\P(\phi)\in \Ker(\iota_\P), \ \forall \P\in\text{Spec}(\mathcal{O}_K)\}
		\end{equation*}
	\end{definition}
	As we have explained above, the weak Mordell-Weil group injects in $S^m(J/K)$ through the map $\delta$. The advantage of working in this smaller group is that we have methods to compute the groups $J(K_\P)/mJ(K_\P)$ for any prime $\P$, so knowing that the restrictions of a 1-cocycle $\phi\in S^m(J/K)$ lie in the image of $\delta_\P$ gives us local conditions that allow the computation of the $m$-Selmer group.
	
	Knowing how to compute the $m$-Selmer group is not enough: in order to produce bounds on the order of the Weak Mordell-Weil group we must also show that it is a finite group. We will do so by explaining how the $m$-Selmer group is contained in a finite cohomology group inside $H^1(G_K,J[m])$, essentially repeating the proof of the weak Mordell-Weil theorem in the language of Galois cohomology. We first translate the property of a finite extension being unramified.
	\begin{definition}
		Let $\P$ be a prime of $K$, let $I_\P\subset G_K$ be an inertia subgroup for $\P$ and let $M$ be a $G_K$-module. We say that a cohomology class $\phi\in H^1(G_K,M)$ is unramified at $\P$ if its restriction to $H^1(I_\P, M)$ is the trivial cohomology class.
		
		If $S$ is a finite set of primes of $K$, we denote by $H^1_S(G_K,M)$ the subgroup of cohomology classes that are unramified at all primes except possibly those in $S$.
	\end{definition} 
	
	
	
	Now we can proof a result similar to \propref{ramifS} which will show that $S^m(J/K)$ is a finite group.
	\begin{proposition}\label{selfinite}Let $K$ be a number field.
		\begin{enumerate}
			\item Let $S$ be a finite set of primes of $K$. If $M$ is a finite $G_K$-module, then the group $H^1_S(G_K,M)$ is finite.
			\item Let $J$ be a Jacobian variety defined over $K$, and $m\ge 2$ an integer. Let $S$ be a set of primes of $K$ containing the archimedean primes of $K$, the primes dividing $m$ and the primes of bad reduction of $J$. Then the $m$-Selmer group $S^m(J/K)$ is contained in $H^1_S(G_K,J[m])$.
		\end{enumerate}
	\end{proposition}
	\begin{remark}
		Since the $m$-torsion of $J$ is a finite $G_K$ module and the set of primes described in part (2) is finite, the two parts of \propref{selfinite} combined imply that the $m$-Selmer group is contained in a finite group so it is finite.
	\end{remark}
	\begin{proof} (1) By definition, the action of $G_K$ on $M$ is continuous when $G_K$ has the Krull topology and $M$ has the discrete topology. Hence, if we fix $m\in M$, the map
			\begin{align*}
				G_K\times \{m\}&\rightarrow M\\
				(\sigma,m)&\mapsto m^\sigma
			\end{align*}
			is continuous. Since $M$ is discrete, the stabilizer of $m$ in $G_K$, denoted by $G_m$, is an open subgroup, because it is the preimage of $m$ by this map. Since $M$ is finite, the set
			\begin{equation*}
				H:=\cap_{m\in M}G_m
			\end{equation*}
			is a subgroup of $G_K$, because it is an intersection of subgroups; and it is open, because it is a finite intersection of open sets. In addition, we have by construction that $m^\sigma=m$ for all $m\in M$ and $\sigma\in H$. 
			
			With the Krull topology on $G_K$, an open subgroup is of the form $\Gal(\overline{K}/L)$, where $L/K$ is a finite extension. Therefore, we have that $G_L$ acts trivially on $M$ (identifying $\Gal(\overline{K}/L))$ with $G_L$). We can now use the restriction-inflation sequence (see \cite{Serre_1979}):
			\begin{equation*}
				\begin{tikzcd}
					0\arrow{r}
					& H^1(G_{L/K},M^{G_L})\arrow{r}
					& H^1(G_K,M)\arrow{r}
					& H^1(G_L,M)
				\end{tikzcd}
			\end{equation*}
			to observe that if we know the result for a finite extension of $K$, then it also holds for $K$. Therefore we may assume, by the above argument, that $G_K$ acts trivially on $M$. In this case, a $1$-cocycle satisfies:
			\begin{equation*}
				\phi(\sigma\tau)=\phi(\sigma)^\tau+\phi(\tau)=\phi(\sigma)+\phi(\tau),
			\end{equation*}
			so cocycles are just group morphisms from $G_K$ to $M$. Moreover, a coboundary has the form:
			\begin{equation*}
				\phi(\sigma)=m^\sigma-m=m-m=0,\text{ for some } m\in M,
			\end{equation*}
			so there are no nontrivial 1-coboundaries. Therefore, we deduce that $H^1(G_K,M)=\Hom(G_K,M)$. Let $m$ be the exponent of $M$ (that is, the smallest positive integer such that every element in $M$ annihilated by $m$). As we saw in the previous chapter, elements in $\Hom(G_K,M)$ correspond to abelian subgroups of finite index in $G_K$ of exponent $m$. By the Fundamental Theorem of Galois Theory, these subgroups correspond to finite abelian extensions of $K$ of exponent $m$. Moreover, elements in $H^1_S(G_K,M)$ correspond to finite abelian extensions of $K$ of exponent $m$ and unramified outside of $S$.
			
			All these extensions lie, by definition, inside the maximal abelian extension of exponent $m$ unramified outside $S$. By \thmref{maxabextension}, this extension is finite so it has a finite number of subextensions, and thus we conclude that $H^1_S(G_K,M)$ is finite.
			
			
			(2) Let $\phi\in S^m(J/K)$ and let $\P$ be a prime not in $S$, which is then non-archimedean by assumption. By definition of the Selmer group, the cohomology class $\phi_\P:=\text{res}_\P(\phi)$ becomes trivial in cohomology when mapped to $H^1(G_\P,J(\overline{K}_\P))[m]$. Therefore, we can find a point $y\in J(\overline{K}_\P)$ such that $\phi_\P(\sigma)=\sigma(y)-y$, for all $\sigma\in G_\P$. Moreover, the point $y$ satisfies the condition $[m]y\in J(K_\P)$. We want to show that for $\sigma\in I_\P$, $\phi_\P(\sigma)$ is the zero cohomology class. Recall that the elements in the inertia group $I_\P$ are those that act trivially modulo $\P$. Therefore, reducing modulo $\P$, we have the equality:
			\begin{equation*}
				\widetilde{\sigma(y)-y}=\widetilde{\sigma(y)}-\widetilde{y}=\widetilde{y}-\widetilde{y}=\widetilde{\mathcal{O}},
			\end{equation*}
			so $\phi(\sigma)$ is zero modulo $\P$. Note also that
			\begin{equation*}
				[m]\phi(\sigma)=([m]y)^\sigma-[m]y=[m]y-[m]y=0,
			\end{equation*}
			because $[m]y$ is invariant by the action of $G_\P$, so $\phi(\sigma)$ is a point of $m$-torsion. By the conditions on $S$, the Jacobian variety $J$ has good reduction at $\P$, and moreover $m$ is coprime to the unique rational prime $p$ lying under $\P$. Therefore, \thmref{torsioninjects} implies that the reduction map
			\begin{equation*}
				J[m]\rightarrow \tilde{J}(\widetilde{k}_\P)
			\end{equation*}
			is injective, so (the restriction to the inertia subgroup of) $\phi(\sigma)$ is equal to zero. In particular it represents the zero class in cohomology so we are done.

	\end{proof}
	Even if we know that the Selmer group is finite, in principle we still have to check an infinite number of conditions to determine whether a cohomology class is in $S^m(J/K)$, one for each prime $\P$ of $K$. Nevertheless, if we choose a set of primes $S$ as in \propref{selfinite}, we have:
	\begin{equation*}
		S^m(J/K)=\{\phi\in H^1_S(G_K,J[m]):\text{res}_\P(\phi)\in \Im(\delta_\P), \ \forall \P\in S\},
	\end{equation*}
	which involves a finite number of computations, by combining Hensel's lemma with estimates on the size of $J(K_\P)/mJ(K_\P)$ that only depend on the dimension of $J$.
	
	We will use the $m$-Selmer group (in fact we will only use $m=2$) to bound the Mordell-Weil rank of some Jacobian varieties, but we first digress for a moment to the question on how could one calculate this rank using the $m$-Selmer group. Since we have an injective map
	\begin{equation*}
		\delta:J(K)/mJ(K)\rightarrow S^m(J/K),
	\end{equation*}
	we are naturally led to ask when is the map $\delta$ an isomorphism, because in that case we could determine the exact rank of $J$ by computing $S^m(J/K)$. To this purpose, we introduce the Tate-Shafarevich group:
	\begin{definition}
		Let $J$ be a Jacobian variety defined over a number field $K$. The Tate-Shafarevich group of $J$ is defined as:
		\begin{equation*}
			\Sh(J/K)=\Ker(H^1(G_K,J(\overline{K}))\rightarrow \prod_\P H^1(G_\P,J(\overline{K}_\P))),
		\end{equation*}
		where the map is the product of the restriction maps in cohomology.
	\end{definition}
	
	By the diagram \eqref{seldiag}, we have that $H^1(G_K,J[m])$ maps surjectively to $\Sh(J/K)[m]$, but in fact it is $S^m(J/K)$ the one that maps surjectively to $\Sh(J/K)[m]$: if $\xi\in\Sh(J/K)[m]$, choose some $\phi\in H^1(G_K,J[m])$ that maps to $\xi$. Then we have, for each prime $\P$:
	\begin{align*}
		\iota_\P\circ\text{res}_\P(\phi)=\text{res}_\P(\xi)=0,
	\end{align*}
	where the first equality follows from the commutativity of \eqref{seldiag}, and the second equality follows from definition of $\Sh(J/K)$. We then have, by exactness of the second row in \eqref{seldiag}, that $\phi\in S^m(J/K)$, as we claimed. This proves that we have the important short exact sequence:
	\begin{equation*}
		\begin{tikzcd}
			0\arrow{r}
			&J(K)/mJ(K)\arrow{r}{\delta}
			&S^m(J/K)\arrow{r}
			&\Sh(J/K)[m]\arrow{r}
			&0
		\end{tikzcd}
	\end{equation*}
	
	Since we know that $S^m(J/K)$ is finite, this implies that $\Sh(J/K)[m]$ is finite. Now we can restate the problem of $\delta$ being an isomorphism as the group $\Sh(J/K)[m]$ being trivial. If the group $\Sh(J/K)$ were finite, then we could find a large enough $m$ such that $\Sh(J/K)[m]=0$, but to this date this is still an open problem:
	\begin{conjecture}(Tate-Shafarevich)
		The group $\Sh(J/K)$ is finite.
	\end{conjecture}
	
	\subsection{An effective algorithm}
	
	Now that we have established the Selmer group as the space that we want to compute, we will explain a series of results given by Schaefer in \cite{Schaefer_1995} to translate our computations into the language of linear algebra over a finite field, and also to replace the cohomology groups with certain finite-dimensional algebras. Specifically, we will work from this point on with the weak Mordell-Weil group for $m=2$, that is, $J(K)/2J(K)$.
	
	To begin with, consider a curve $C$ defined over a field $K$ of characteristic 0 by the equation $y^2=f(x)$, where $f(x)$ is a monic, separable polynomial of odd degree $2g+1$ and coefficients in $K$, thus the curve $C$ has genus $g$. Fix an algebraic closure of $K$, denoted by $\overline{K}$. Since we want to construct the Jacobian of $C$, we must ensure that $C$ is non-singular. The curve $C$ is only singular at the (unique) point at infinity, so we solve this singularity by taking the normalization of the curve. Explicitly, what we are doing at this point $\infty$ is to take the integral closure of the local ring $K[C]_\infty$. By abuse of notation, we call this normalization again $C$, and consider its Jacobian $J:=\Jac(C)$.
	
	Since we are fixing $\infty$ as the base point, the embedding of $C$ into $J$ is given by:
	\begin{align*}
		\iota:C&\hookrightarrow J\\
		P&\mapsto [P-\infty]
	\end{align*}
	
	Let $\{\alpha_1\ldots,\alpha_d\}$ be the $d$ distinct roots of $f(x)$ in $\overline{K}$. Then the divisors including $[(\alpha_i,0)-\infty]$, for $i=1,\ldots,d-1$ form a basis for $J[2]\cong (\Zmod{2})^{d-1}$. The remaining divisor $[(\alpha_d,0)-\infty]$ satisfies:
	\begin{equation*}
		[(\alpha_d,0)-\infty]=\sum_{i=1}^{d-1}[(\alpha_i,0)-\infty],
	\end{equation*}
	because
	\begin{equation*}
		\Div(y)=\sum_{i=1}^{d}[(\alpha_i,0)-\infty]
	\end{equation*}
	(and these divisors are points of 2-torsion, so the sign does not matter).
	
	We now define the algebra $L=K[T]/(f(T))$ over $K$, which is étale because, if we define $\overline{L}:=L\otimes \overline{K}$ we have, by the Chinese remainder theorem:
	\begin{equation*}
		\overline{L}\cong\overline{K}[T]/(f(T))\cong\overline{K}[T]/(T-\alpha_1)\oplus\overset{(d}{\ldots}\oplus\overline{K}[T]/(T-\alpha_d)\cong \overline{K}\oplus \overset{(d}{\ldots} \oplus\overline{K}.
	\end{equation*}
	
	We call $\varphi$ the composition of all these isomorphisms. We can now let $G_K$ act trivially on $T$, so that $\overline{L}$ becomes a $G_K$-module extending this action to $\overline{K}[T]$ (the action descends to the quotient because $f(T)\in K[T]$). It is important to study this Galois action on $\overline{K}^d$ when we consider the isomorphism of algebras above. We transport the structure, that is, we define the action of $G_K$ in $\overline{K}^d$ by stating that $\varphi$ is a $G_K$-module morphism. If $h(T)\in\overline{K}[T]$, then the image of $[h(T)]$ in $\overline{K}^d$ is the $d$-tuple $(h(\alpha_1),\ldots,h(\alpha_d))$. Therefore, if we want to find the value of $(\lambda_1,\ldots,\lambda_d)^\sigma$, where $\lambda_i\in\overline{K}$ and $\sigma\in G_K$, we must first find the preimage of the tuple. We do so by using the Lagrange interpolation polynomial:
	\begin{equation*}
		h(T)=\sum_{j=1}^d\lambda_j l_j(T),\text{ with } l_j(T)=\prod_{m\neq j}\frac{T-\alpha_m}{\alpha_j-\alpha_m},
	\end{equation*}
	satisfies that $h(\alpha_i)=\lambda_i$ for each $1\le i \le d$, because $l_j(\alpha_i)=\delta_{ij}$. We have then found that the class of $h(T)$ maps to $(\lambda_1,\ldots,\lambda_d)$. By the numbering that we have fixed on the roots $\alpha_i$, each $\sigma\in G_K$ defines a certain permutation of the symmetric group $\Sigma_d$ that we write as $\alpha_i^\sigma=\alpha_{\sigma(i)}$. Then the action of $\sigma$ at $l_j(T)$ is:
	\begin{equation*}
		l_j^\sigma(T)=\prod_{m\neq j}\frac{T-\alpha_{\sigma(m)}}{\alpha_{\sigma(j)}-\alpha_{\sigma(m)}}=l_{\sigma(j)}(T),
	\end{equation*}
	so
	\begin{equation*}
		h^\sigma(T)=\sum_{j=1}^d \lambda_j^\sigma l_{\sigma(j)}(T).
	\end{equation*}
	Therefore the action on $(\lambda_1,\ldots,\lambda_d)$ is:
	\begin{equation*}
		(\lambda_1,\ldots,\lambda_d)^\sigma=\varphi([h(T)])^\sigma=\varphi([h^\sigma(T)])=(h^\sigma(\alpha_1),\ldots,h^\sigma(\alpha_d)).
	\end{equation*}
	Finally, observe that $l_{\sigma(j)}(\alpha_{\sigma(i)})=\delta_{ij}$ because $\sigma$ is bijective, so $l_{\sigma(j)}(\alpha_i)=l_{\sigma(j)}(\alpha_{\sigma\sigma^{-1}(i)})=\delta_{\sigma^{-1}(i)j}$.
	
	Thus we obtain:
	\begin{equation*}
		(\lambda_1,\ldots,\lambda_d)^\sigma=(\lambda_{\sigma^{-1}(1)}^\sigma,\ldots,\lambda_{\sigma^{-1}(d)}^\sigma).
	\end{equation*}
	
	In particular we see that the action on $\overline{K}^d$ permutes each $\overline{K}$. We can think of $L$ as the elements in $\overline{L}$ that are invariant by this action of $G_K$. Whether we talk about tuples of points, divisors or functions we will have fixed the same permutation of the Galois action. For instance, we consider the $d$-tuple of divisors $([(\alpha_1,0)-\infty],\ldots,[(\alpha_d,0)-\infty])$, which we call $P_T$. Applying any $\sigma\in G_K$ to this tuple, we get:
	\begin{align*}
		([(\alpha_1,0)-\infty],\ldots,[(\alpha_d,0)-\infty])^\sigma&=([(\alpha_{\sigma^{-1}(1)},0)-\infty]^\sigma,\ldots,[(\alpha_{\sigma^{-1}(d)},0)-\infty]^\sigma)\\
		&=([(\alpha_1,0)-\infty],\ldots,[(\alpha_d,0)-\infty])
	\end{align*}
	so it is a $G_K$ invariant tuple. In addition, each divisor $[(\alpha_i,0)-\infty]$ is of 2-torsion in $J$. Consider the square roots of unity in $\overline{L}$, which we denote by $\mu_2(\overline{L})$. By the isomorphism $\varphi$, we have that $\mu_2(\overline{L})$ is isomorphic to $(\pm 1)^d$. We now recall the definition and basic properties of the Weil pairing in the following proposition.
	\begin{proposition}\label{Weilpairing}
		Let $J=\Jac(C)$ for some smooth projective curve $C$ of genus $g\ge 0$. Let $T\in J[2]$ be a 2-torsion point. Then there exists a function $f\in\overline{K}(C)$ such that:
		\begin{equation*}
			\Div(f)=2T-2\infty.
		\end{equation*}
		Moreover,
		
		
		
		The pairing
		\begin{align*}
			J[m]\times J[m]&\rightarrow \mu_2(\overline{K})=\{\pm 1\}\\
			(T,S)&\mapsto e_2(T,S)
		\end{align*}
		called the Weil pairing, has the following properties:
		\begin{enumerate}
			\item a
		\end{enumerate}
	\end{proposition}
		
	\begin{remark}
		The Weil pairing can be defined on the $m$-torsion of $J$ for any $m$. The property that the pairing is symmetric is exclusive to $m=2$.
	\end{remark}
	For a proof of all these claims, we refer the reader to .
	
	We can then define a map
	\begin{gather*}
		w:J[2]	\rightarrow\mu_2(\overline{L})\\
		P\mapsto (e_2(P,[(\alpha_1,0)-\infty]),\ldots, e_2(P,[(\alpha_d,0)-\infty]))
	\end{gather*}
	which is a homomorphism of groups by bilinearity of the Weil pairing. Moreover, it is a morphism of $G_K$ modules: if $P\in J[2]$ and $\sigma\in G_K$,
	\begin{align*}
		w(P^\sigma)&=(e_2(P^\sigma,[(\alpha_1,0)-\infty]),\ldots, e_2(P^\sigma,[(\alpha_d,0)-\infty]))\\
		&=(e_2(P^\sigma,[(\alpha^\sigma_{\sigma^{-1}(1)},0)-\infty]),\ldots, e_2(P^\sigma,[(\alpha^\sigma_{\sigma^{-1}(d)},0)-\infty]))\\
		&=(e_2(P,[(\alpha_{\sigma^{-1}(1)},0)-\infty])^\sigma,\ldots, e_2(P,[(\alpha_{\sigma^{-1}(d)},0)-\infty])^\sigma)\quad\text{(By $G_K$ invariance of $e_2$)}\\
		&=w(P)^\sigma\quad\text{(By definition)}.
	\end{align*}
By functoriality of Galois cohomology, the morphism $w$ induces a group homomorphism
\begin{equation*}
	H^1(w):H^1(G_K,J[2])\rightarrow H^1(G_K,\mu_2(\overline{L}))
\end{equation*}
which, by abuse of notation, we denote again by $w$ (it will be clear from the context which one we are using).

Observe that we have the following exact sequence of $G_K$-modules:
\begin{equation*}
	\begin{tikzcd}
		0\arrow{r}
		&\mu_2(\overline{K})\arrow[hook]{r}
		&\overline{K}^*\arrow{r}{z\mapsto z^2}
		&\overline{K}^*\arrow{r}
		&0,
	\end{tikzcd}
\end{equation*}
where the action on $\mu_2(\overline{K})=\{\pm1\}$ is trivial. Then by taking $G_K$ invariants we obtain the long exact sequence in cohomology:
\begin{equation*}
	\begin{tikzcd}
		0\arrow{r}
		&\mu_2(K)\arrow{r}
		&K^*\arrow{r}{z\mapsto z^2}
		&K^*\arrow{r}{\delta}
		&\phantom{ }
		\\ \arrow{r}{\delta}
		&H^1(G_K,\mu_2(\overline{K}))\arrow{r}
		&H^1(G_K,\overline{K}^*)\arrow{r}
		& \ldots,
	\end{tikzcd}
\end{equation*}
	We remark that this situation is completely analogous to the one in \eqref{kummersequence}. By Hilbert's Theorem 90, we have $H^1(G_K,\overline{K}^*)=\{0\}$ (it follows from the case when $G$ is finite because $G_K$ is constructed as a projective limit of finite groups, which commutes with Galois cohomology). Therefore, the map $\delta$ is surjective. The kernel of this map is, by exactness, the image of the map $z\mapsto z^2$, which is by definition $(K^*)^2$, so we obtain an isomorphism between $K^*/(K^*)^2$ and $H^1(G_K,\mu_2(\overline{K}))$. Since $\overline{L}$ is a finite dimensional $\overline{K}$-algebra (of dimension $d$) and it has a unit (the tuple with 1 in each component), we deduce that $H^1(G_K,\overline{L}^*)=\{0\}$ (see \cite{Serre_1979}, page 152). By the same reasoning as above, we obtain an isomorphism between $H^1(G_K,\mu_2(\overline{L}))$ and $L^*/(L^*)^2$, which we denote by $\kappa$. This isomorphism is the inverse of the map that sends $a\in L^*/(L^*)^2$ to the 1-cocycle represented by $\sigma\mapsto b^\sigma/b$, where $b\in\overline{L}$ is such that $b^2=a$.
	
	We have to define one more map, coming from the norm function:
	\begin{align*}
		N:\overline{L}&\rightarrow\overline{K}\\
		(\lambda_1,\ldots,\lambda_d)&\mapsto \lambda_1\cdot\ldots\cdot\lambda_d
	\end{align*}
	which clearly maps $\overline{L}^*$ to $\overline{K}^*$. If $\lambda=(\lambda_1,\ldots,\lambda_d)$ is in $\overline{L}$, then we find that, for any $\sigma\in G_K$:
	\begin{align*}
		N(\lambda^\sigma)&=N(\lambda_{\sigma^{-1}(1)}^\sigma,\ldots,\lambda_{\sigma^{-1}(d)}^\sigma)\\
		&=\lambda_{\sigma^{-1}(1)}^\sigma\cdot\ldots\cdot\lambda_{\sigma^{-1}(d)}^\sigma\\
		&=\lambda_1^\sigma\cdot\ldots\cdot\lambda_d^\sigma\quad(\sigma\text{ is bijective})\\
		&=(\lambda_1\cdot\ldots\cdot\lambda_d)^\sigma=N(\lambda)^\sigma
	\end{align*}
	so $N$ is a morphism of $G_K$-modules. In particular, if $\lambda$ is in $L$, so it is invariant by $G_K$, then $N(\lambda)$ is invariant by $G_K$ so it is in $K$. Therefore the norm descends to a map $N:L^*\rightarrow K^*$. If we now compose this map with the projection $K^*\rightarrow K^*/(K^*)^2$, we see that $(L^*)^2$ is in the kernel of the composition so finally obtain a map:
	\begin{equation*}
		L^*/(L^*)^2\rightarrow K^*/(K^*)^2
	\end{equation*}
	that we call again $N$. Similarly, the norm $\overline{L}\rightarrow \overline{K}$ restricts well to the second roots of unity $\mu_2(\overline{L})\rightarrow\mu_2(\overline{K})$, so we get an induced map in cohomology $H^1(G_K,\mu_2(\overline{L}))\rightarrow H^1(G_K,\mu_2(\overline{K}))$. We will call all these maps $N$, but it will be clear from the context which one we are using in any case.
	
	The key point now is that it is easy to work in the space $L^*/(L^*)^2$ and moreover in the kernel of $N$, because $d-1$ components of a tuple determine the remaining one. The following theorem allows us to replace the cohomology group $H^1(G_K,J[2])$ (so in particular the $2$-Selmer group) with this nice space.
	\begin{theorem}\label{kwisomorphism}
		The composition of the maps $\kappa\circ w$ is an isomorphism between $H^1(G_K),J[2])$ and the kernel of the norm $N:L^*/(L^*)^2\rightarrow K^*/(K^*)^2$.
	\end{theorem}
	\begin{proof}
		We claim that
		\begin{equation*}
			\begin{tikzcd}
				0\arrow{r}
				&J[2]\arrow{r}{w}
				&\mu_2(\overline{L})\arrow{r}{N}
				&\mu_2(\overline{K})\arrow{r}
				&1
			\end{tikzcd}
		\end{equation*}
		is an exact sequence of $G_K$-modules, the action on $J[2]$, $\mu_2(\overline{L})$ and $\mu_2(\overline{K})$ being the same one as on $J(\overline{K})$, $\overline{L}$ and $\overline{K}$, respectively, as subsets. We have already seen that $w$ and $N$ are morphisms of $G_K$-modules. First, we prove that $w$ is injective. Let $P\in J[2]$ be a 2-torsion point such that $w(P)=1$. By definition, we have that $e_2(P,[(\alpha_j,0)]-\infty)=1$, for all $j$. Since the divisors $[(\alpha_j,0)]-\infty$ form a basis for $J[2]$ and $e_2$ is bilinear, we find that $e_2(P,T)=1$ for all $T\in J[2]$. Since $e_2$ is non-degenerate, this implies that $P=0$.
		
		The norm is clearly surjective, since the tuple $(-1,1,\ldots,1)$ maps to $-1$.
		
		Next, we see that the image of $w$ is contained in the kernel of $N$. Let $P$ be a point of 2-torsion. Then, we compute:
		\begin{align*}
			N(w(P))&=\prod_{i=1}^d e_2(P,[(\alpha_i,0)]-\infty)\\
			&=e_2\left(P,\sum_{i=1}^d [(\alpha_i,0)]-\infty\right)    \quad (e_2 \text{ is bilinear})\\
			&=e_2(P,0)=1,
		\end{align*}
		which shows what we wanted. To prove the reverse inclusion, we notice that each term in the sequence is a finitely dimensional vector space over $\mathbb{F}_2$. Concretely, $J[2]$ has $2$-rank $2g=d-1$, $\mu_2(\overline{L})$ has 2-rank $d$ and $\mu_2(\overline{K})$ has 2-rank 1. By the Rank-Nullity theorem, we have:
		\begin{equation*}
			\dim_2(\Im(w))=\dim_2(J[2])-\dim_2(\Ker(w))=d-1 \quad (w\text{ is injective})
		\end{equation*}
		and
		\begin{equation*}		
			\dim_2(\Ker(N))=\dim_2(\mu_2(\overline{L}))-\dim_2(\Im(N))=d-1 \quad (N\text{ is surjective}),
		\end{equation*}
		so we must have $\Im(w)=\Ker(N)$. Now we deduce that there is a long exact sequence in cohomology by taking $G_K$ invariants:
		\begin{equation*}
			\begin{tikzcd}
				0\arrow{r}
				&J(K)[2]\arrow{r}
				&\mu_2(L)\arrow{r}
				&\mu_2(K)\arrow{r}
				&\phantom{ }
				\\ \arrow{r}
				&H^1(G_K,J[2])\arrow{r}{w}
				&H^1(G_K,\mu_2(\overline{L}))\arrow{r}{N}
				&H^1(G_K,\mu_2(\overline{K}))\arrow{r}
				& \ldots
			\end{tikzcd}
		\end{equation*}
		The tuple $(-1,\ldots,-1)\in\mu_2(\overline{L})$ is clearly $G_K$ invariant, so it lies in $\mu_2(L)$. Its norm is $(-1)^d$ and, since $d$ is odd, it is $-1$ so the norm $\mu_2(L)\rightarrow\mu_2(K)$ is surjective. Therefore, the induced map $w:H^1(G_K,J[2])\rightarrow H^1(G_K,\mu_2(\overline{L}))$ is injective, so it is an isomorphism between $H^1(G_K,J[2])$ and $\Im(w)=\Ker(N)$. But we have the following commutative diagram:
		\begin{equation*}
		\begin{tikzcd}
			1\arrow{r}
			&H^1(G_K,\mu_2(\overline{L}))\arrow{r}{\kappa}\arrow{d}{N}
			&L^*/(L^*)^2\arrow{r}\arrow{d}{N}
			&1\\
			1\arrow{r}
			&H^1(G_K,\mu_2(\overline{K}))\arrow{r}
			&K^*/(K^*)^2\arrow{r}
			&1
		\end{tikzcd}
		\end{equation*}
		so we can replace $H^1(G_K,\mu_2(\overline{L}))$ by $L^*/(L^*)^2$ using $\kappa$ and deduce that $\kappa\circ w$ is an isomorphism between $H^1(G_K,J[2])$ and the kernel of the norm from $L^*/(L^*)^2$ to $K^*/(K^*)^2$, as we wanted.
	\end{proof}
	
	The next step in our work is to replace the map $\delta:J(K)/2J(K)\rightarrow H^1(G_K,J[2])$ from the previous section by a suitable map that provides the same information but is easier to compute.
	
	In general, let $f$ be a function on $C$ and $R=\sum n_iR_i$ be a divisor whose support is disjoint from $\Div(f)$. Then we can define $f(R)$ to be $\prod f(R_i)^{n_i}$, which is an element in $\overline{K}^*$ ($f$ has no poles or zeros among the $R_i$).
	
	In particular, consider the function $(x-\alpha_j)$, whose divisor has support $(\alpha_j,0)$ and $\infty$. Let $R=\sum n_i R_i$ be a divisor on $C$ of degree 0 which is defined over $K$. Recall that $R$ is defined over $K$ if, for all $\sigma\in G_K$, we have $R^\sigma=\sum n_iR_i^\sigma=R$, which is a weaker assertion than requiring each $R_i$ to be a point in $C(K)$. Assume further that none of the points $R_i$ is a Weierstrass point, i.e., none of the $R_i$ are the point at infinity nor they have $0$ as $y$-coordinate. Then the support of $R$ is disjoint from $\Div(x-\alpha_j)$, for each $j$. We can therefore define, for each $j$:
	\begin{equation*}
		(x-\alpha_j)(R)=\prod(x(R_i)-\alpha_j)^{n_i},
	\end{equation*}
	which is an element in $\overline{K}^*$. Since the isomorphism $\varphi$ maps $T$ to the tuple $(\alpha_1,\ldots,\alpha_d)$, we can define a function $(x-T)$ only on these kind of divisors with values on $\overline{L}^*$ as:
	\begin{equation*}
		(x-T)(R)=\prod (x(R_i)-T)^{n_i}
	\end{equation*}
	which maps to $(x-\alpha_j)(R)$ in each $\overline{K}^*$ through $\varphi$.
	
	We would like to define a function on all $J(K)$ using $(x-T)$. The first step to do that is to show that the map $(x-T)$ does not depend on the linear class of divisors (with our particular restriction on their support) that we choose, up to squares. We will use the following result known as Weil reciprocity:
	\begin{proposition}(Weil reciprocity)
		Let $f$ and $g$ be two rational functions on a smooth complete curve $C$ with disjoint support. Then
		\begin{equation*}
			f(\Div(g))=g(\Div(f)).
		\end{equation*}
	\end{proposition}
	For a proof of Weil reciprocity, we refer to \cite{Arbarello_Cornalba_Griffiths_Harris_1985}.
	\begin{lemma}\label{x-Twelldefined}
		The function $x-T$ is a well defined map from points of $J(K)$ represented by divisors of degree 0, defined over $K$ with support disjoint from the Weierstrass points, to $L^*/(L^*)^2$.
	\end{lemma}
	\begin{proof}
		Let $D_1,D_2$ be two linearly equivalent divisors of degree 0, defined over $K$, and with support disjoint from the Weierstrass points. We have to show that their values at $x-T$ are in $L^*$ and differ only by a square.
		
		Choose a function $h$ defined over $K$ such that $\Div(h)=D_1-D_2$. Since then $\Div(h)$ has support disjoint from the Weierstrass points, we have:
		\begin{align*}
			(x-\alpha_j)(D_1-D_2)&=(x-\alpha_j)(\Div(h))\\
			&=h(\Div(x-\alpha_j))\quad\text{(Weil reciprocity)}\\
			&=h(2(\alpha_j,0)-2\infty)\\
			&=h((\alpha_j,0)-\infty)^2.
		\end{align*}
		By the above calculation, we only know that $(x-\alpha_j)(D_1-D_2)\in\overline{K}^*$, but we can collect this information for each $j$ to conclude the following equality of $d$-tuples in $\overline{L}^*$:
		\begin{equation*}
			(x-T)(D_1-D_2)=h(P_T)^2.
		\end{equation*}
		Since $P_T$ is fixed by the Galois action and $h$ is defined over $K$, we find that $h(P_T)$ is in $L^*$, so $(x-T)(D_1-D_2)$ is in $(L^*)^2$, as we wanted.
	\end{proof}
	The second step in our study of $x-T$ is to show that we can always find a divisor with our assumptions inside any linear equivalence class in $J(K)$.
	\begin{lemma}\label{x-Tdeformation}
		Any point in $J(K)$ can be represented by a divisor whose support is disjoint from the Weierstrass points. In addition, the map $(x-T)$ can be computed as follows:
		Let $P$ be a point in $J(K)$ that is represented by a divisor which includes in its formal sum expression a divisor of the form:
		\begin{equation*}
			D=\left(\sum_{i=1}^r Q^{\sigma_i}\right) -r\infty,
		\end{equation*}
		where $Q$ is a point of $C$ with $y$-coordinate distinct from zero and the $Q^{\sigma_i}$ are the $r$ distinct conjugates of $Q$ over $K$. Then:
		\begin{equation*}
			(x-T)(P)\equiv \prod_{i=1}^r(x(Q^{\sigma_i})-T)\pmod{(L^*)^2}.
		\end{equation*}
		Let $P$ be a point in $J(K)$ that is represented by a divisor which includes in its formal sum expression a divisor of the form:
		\begin{equation*}
			D=\left(\sum_{i=1}^r [(\alpha_i,0)]\right) -r\infty,
		\end{equation*}
		where the $\alpha_i$ are roots of $f(T)$ which are conjugate over $K$. Denote also by $\alpha_i$, for $r+1\le i\le d$ the remaining roots of $f(T)$, if there are any. Then:
		\begin{equation*}
			(x-T)(P)\equiv \prod_{i=1}^r (\alpha_i-T)+\prod_{i=r+1}^d (\alpha_i-T)\pmod{(L^*)^2}.
		\end{equation*}
	\end{lemma}
	\begin{proof}
		Consider a point in $J(K)$ represented by a divisor whose support contains some Weierstrass point. We will find a linearly equivalent divisor (so it represents the same point in $J(K)$) whose support does not contain any Weierstrass points. Fix a point $Q_1$ in the support of this divisor distinct from $\infty$, and let $Q_1^{\sigma_1},\ldots,Q_1^{\sigma_r}$ be the $r$ distinct conjugates of $Q_1$ over $K$, with $\sigma_i\in G_K$. Then the divisor
		\begin{equation*}
			D=\left(\sum_{i=1}^r Q_1^{\sigma_i}\right) -r\infty
		\end{equation*}
		has degree 0 and is $G_K$ invariant (i.e., it is defined over $K$). Possibly selecting differents sets $\{\sigma_i\}$, any degree 0 divisor defined over $K$ can be written as sums and differences of divisors of the form of $D$, so it is enough to prove the lemma for those divisors. We fix some notation: let $Q_1=(x_1,y_1)$ be a point in $C$ and let $\{\sigma_i\}$ be a set of $r$ morphisms such that $Q_1^{\sigma_1},\ldots,Q_1^{\sigma_r}$ are the conjugates of $Q_1$ over $K$. Then for any point $Q$ of $C$ defined over $K(x_1,y_1)$ we define:
		\begin{equation*}
			\overline{Q}=\sum_{i=1}^r Q^{\sigma_i}.
		\end{equation*}
		In particular, we have that $D=\overline{Q_1}-r\infty$. We will distinguish between two cases depending on $Q_1$.
		
		If $y_1$ is not zero, then we have to get rid of the point at infinity in $D$. Counted with multiplicities, the equation $y_1^2=f(T)$ has $d$ solutions, that correspond to points $Q_1,\ldots,Q_d$ with $y$-coordinate $y_1$. Let $R_1=(x_1,-y_1)$. Then we have the following identities:
		\begin{equation*}
			\Div(y-y_1^{\sigma_i})=\left(\sum_{j=1}^d Q_j^{\sigma_i}\right)-d\infty
		\end{equation*}
		and
		\begin{equation*}
			\Div(x-x_1^{\sigma_i})= Q_1^{\sigma_i}+R_1^{\sigma_i}-2\infty.
		\end{equation*}
		Therefore with the notation fixed above, we have:
		\begin{equation*}
				\Div\left(\prod_{i=1}^r \frac{1}{y-y_1^{\sigma_i}}\right)=dr\infty-\overline{Q_1}-\ldots-\overline{Q_d}
		\end{equation*}
		and
		\begin{equation*}
			\Div\left(\prod_{i=1}^r (x-x_1^{\sigma_i})^g\right)=g\overline{Q_1}+g\overline{R_1}-2gr\infty.
		\end{equation*}
		and these two functions are defined over $K$. By adding their divisors to $D$, we have:
		\begin{align*}
			D=\overline{Q_1}-r\infty&\sim \overline{Q_1}-r\infty+dr\infty-\overline{Q_1}-\ldots-\overline{Q_d}+g\overline{Q_1}+g\overline{R_1}-2gr\infty\\
			&=g\overline{Q_1}+g\overline{R_1}-\overline{Q_2}-\ldots-\overline{Q_d},
		\end{align*}
		because $d=2g+1$. Note that this divisor, which is equivalent to $D$, does not have Weierstrass points in its support. Therefore, if $D$ represents the point $P\in J(K)$, we can compute $(x-T)(P)$ by using the equivalent divisor that we have found and the definition of $x-T$:
		\begin{align*}
			(x-T)(P)&\equiv \prod_{i=1}^r \frac{(x(Q_1^{\sigma_i})-T)^g(x(R_1^{\sigma_i})-T)^g}{(x(Q_2^{\sigma_i})-T)\cdot\ldots\cdot(x(Q_d^{\sigma_i})-T)}\\
			&\equiv \prod_{i=1}^r \frac{(x(Q_1^{\sigma_i})-T)^{2g}}{(x(Q_2^{\sigma_i})-T)\cdot\ldots\cdot(x(Q_d^{\sigma_i})-T)}\quad (x(Q_1)=x(R_1))\\
			&\equiv \prod_{i=1}^r \frac{(x(Q_1^{\sigma_i})-T)^d}{(x(Q_1^{\sigma_i})-T)\cdot\ldots\cdot(x(Q_d^{\sigma_i})-T)}\quad (2g=d-1)\\
			&\equiv \prod_{i=1}^r \frac{(x(Q_1^{\sigma_i})-T)^d}{(y(Q_1^{\sigma_i}))^2-f(T)}\\
			&\equiv \prod_{i=1}^r \frac{(x(Q_1^{\sigma_i})-T)(x(Q_1^{\sigma_i})-T)^{2g}}{(y(Q_1^{\sigma_i}))^2}\quad (f(T)\text{ is zero in } L)\\
			&\equiv \prod_{i=1}^r (x(Q_1^{\sigma_i})-T)\pmod{(L^*)^2}.
		\end{align*}
	Now we deal with the case $y_1=0$. We have $Q_1=(\alpha_1,0)$, where $\alpha_1$ is one of the roots of $f(T)$. If $\alpha_1$ has $d$ conjugates over $K$, then $D=\overline{Q_1}-d\infty=\Div(y)$, which is zero. We can thus assume that $r$, the number of conjugates of $\alpha_1$ over $K$, is strictly less than $d$. Observe that the function $\prod_{i=r+1}^d (x-\alpha_i)/y$ has divisor:
	\begin{align*}
		\Div\left(\prod_{i=r+1}^d(x-\alpha_i)/y\right)&=\sum_{i=r+1}^d 2(\alpha_i,0)-2(d-r)\infty-\sum_{i=1}^d(\alpha_i,0)+d\infty\\
		&=\sum_{i=r+1}^d(\alpha_i,0)-\sum_{i=1}^r(\alpha_i,0)+(-d+2r)\infty.
	\end{align*}
	If we now add this divisor to our divisor $D=(\alpha_1,0)+\ldots+(\alpha_r,0)-r\infty$, we obtain:
	\begin{align*}
		D&\sim \sum_{i=1}^r(\alpha_i,0)-r\infty+\sum_{i=r+1}^d(\alpha_i,0)-\sum_{i=1}^r(\alpha_i,0)+(-d+2r)\infty\\
		&=\sum_{r+1}^d(\alpha_i,0)-(d-r)\infty.
	\end{align*}
	
	
	Therefore this sum, which does not affect the linear class of $D$, swaps between the $r$ first roots $\alpha_i$ and the last $d-r$. Since $d$ is odd, one of these quantities is strictly smaller than $d/2$. Therefore, by adding that divisor if necessary, we can assume that $r<d/2$.
	Reindexing if necessary, we have that $\alpha_1,\ldots,\alpha_r$ are the $r$ conjugates of $\alpha_1$. Consider the function $g=y-\prod_{i=1}^r(x-\alpha_i)$, so we have:
	\begin{equation*}
		\Div(g)=(\alpha_1,0)+\ldots+(\alpha_r,0)+\sum_{i=r+1}^d P_i-d\infty,
	\end{equation*}
	where $P_i=(x_i,y_i)$ are points with $y_i\neq 0$ and $x_i$ are the roots of the polynomial:
	\begin{equation*}
		\prod_{i=r+1}^d(x-\alpha_i)-\prod_{i=1}^r(x-\alpha_i).
	\end{equation*}
	(Note that our assumption $r<d/2$ implies that this polynomial has degree $d-r$). The presence of the points $P_i$ in the divisor of $g$ is explained as follows: if $x_i$ is as above, then the equation $y^2=f(x_i)$ has two solutions corresponding to the points $(x_i,\pm y_i)$. Then we have:
	\begin{align*}
		\left(y_i-\prod_{i=1}^r(x_i-\alpha_i)\right)&\cdot\left(-y_i-\prod_{i=1}^r(x_i-\alpha_i)\right)=\prod_{i=1}^r(x_i-\alpha_i)\prod_{i=1}^r(x_i-\alpha_i)-y_i^2\\
		&=\prod_{i=1}^r(x_i-\alpha_i)\prod_{i=r+1}^d(x_i-\alpha_i)-y_i^2\quad(\text{By definition of the }x_i)\\
		&=\prod_{i=1}^d(x_i-\alpha_i)-y_i^2=f(x_i)-y_i^2=0,
	\end{align*}
	so exactly one of the points $(x_i,\pm y_i)$ is a zero of the function $g$. Choosing the correct one for each $i$ and renaming it to $(x_i,y_i)$ yields the result.
	
	Define $R_i=(x_i,-y_i)$ and $h=\prod_{i=r+1}^d (x-x_i)$, so we have:
	\begin{equation*}
		\Div(h)=\sum_{i=r+1}^d P_i+\sum_{i=r+1}^d R_i -2(d-r)\infty.
	\end{equation*}
	We can add the divisor $\Div(h)$ and substract $\Div(g)$ to $D$ and obtain:
	\begin{align*}
		D=\overline{Q_1}-r\infty&\sim \overline{Q_1}-r\infty+\Div(h)-\Div(g)\\
		&=\sum_{i=r+1}^d R_i-(d-r)\infty,
	\end{align*}
	and we are in the same case as before. Therefore, we have:
	\begin{align*}
		(x-T)(P)&\equiv \prod_{i=r+1}^d(x_i-T)\equiv (-1)^{d-r}\prod_{i=r+1}^d(T-x_i)\\
		&\equiv(-1)^{d-r}\left(\prod_{i=r+1}^d(T-\alpha_i)-\prod_{i=1}^r(T-\alpha_i)\right)       \quad(\text{By definition of the }x_i)\\
		&\equiv \prod_{i=r+1}^d(\alpha_i-T)-\prod_{i=1}^r(\alpha_i-T) \pmod{(L^*)^2},
	\end{align*}
	where the last equality follows from the fact that $d$ is odd.
	\end{proof}
	
	The formula to compute $x-T$ is now the following: for any point in $J(K)$, we can find a divisor representing it of the desired form, by \lemref{x-Tdeformation}. The choice of such a divisor is not important by \lemref{x-Twelldefined}. Then the value of $x-T$ is given again in \lemref{x-Tdeformation}.
	
	Now that we have obtained a well defined function from $J(K)$ to $L^*/(L^*)^2$, observe the following: if $P\in 2J(K)$, find $Q\in J(K)$ such that $P=2Q$. By \lemref{x-Tdeformation}, we can assume that $P$ has disjoint support from the Weierstrass points (thus the same holds for $Q$) so we can compute:
	\begin{equation*}
		(x-T)(P)=(x-T)(2Q)=((x-T)(Q))^2=1\pmod {(L^*)^2},
	\end{equation*}
	and similarly a point $P$ such that $(x-T)(P)=1$ in $L^*/(L^*)^2$ must be of the form $2Q$ for some $Q\in J(K)$.
	
	Therefore, the map $x-T$ descends to an injective map from $J(K)/2J(K)$ to $L^*/(L^*)^2$. Thanks to the next theorem, this is the function that we will use to compute the image of $J(K)/2J(K)$ inside the Selmer group.
	\begin{theorem}\label{x-Tinjection}
		The maps $x-T$ and $\kappa\circ w\circ \delta$ are identical as injections from $J(K)/2J(K)$ to $L^*/(L^*)^2$.
	\end{theorem}
	\begin{proof}
		Let $P$ be a degree 0 divisor defined over $K$ and whose support is disjoint from the Weierstrass points. By \lemref{x-Tdeformation}, such a choice of $P$ can represent any point in $J(K)$, and by \lemref{x-Twelldefined} the precise choice of $P$ is not important. We will compute first the value of $\kappa\circ w\circ \delta(P)$. To compute $\delta(P$, we must choose a divisor $Q$ (defined over $\overline{K})$ such that $2Q=P$. Also by \lemref{x-Tdeformation}, we can choose $Q$ with support disjoint from the Weierstrass points. Then by definition, $\delta(P)$ is the class of the cocycle $\sigma\mapsto Q^\sigma-Q$. Then $w\circ\delta(P)$ is the class of the cocycle $\sigma\mapsto w(Q^\sigma-Q)$. We thus have to compute
		\begin{equation*}
			e_2(Q^\sigma-Q,[(\alpha_j,0)]-\infty),
		\end{equation*}
		for all $j$. By definition of the Weil pairing, we have two find functions $h_1,h_2$ such that
		\begin{equation*}
			\Div(h_1)=2Q^\sigma-2Q
		\end{equation*}
		and
		\begin{equation*}
			\Div(h_2)=2[(\alpha_j,0)]-2\infty.
		\end{equation*}
		We know that there is a function $g$ on $C$ such that $\Div(g)=2Q-P$. Then we have that $\Div(g^\sigma)=2Q^\sigma-P$, so $\Div(g^\sigma/g)=2Q^\sigma-P-2Q+P=2Q^\sigma-2Q$. Also recall that $\Div(x-\alpha_j)=2[(\alpha_j,0)]-2\infty$. Then we have that
		\begin{equation*}
			e_2(Q^\sigma-Q,[(\alpha_j,0)]-\infty)=\frac{(x-\alpha_j)(Q^\sigma-Q)}{(g^\sigma/g)([(\alpha_j,0)]-\infty)}.
		\end{equation*}
		We can collect these equalities for each $j$ in the following equality of $d$-tuples:
		\begin{equation}
			e_2(Q^\sigma-Q,P_T)=\frac{(x-T)(Q^\sigma-Q)}{(g^\sigma/g)(P_T)}=\frac{\beta^\sigma}{\beta},
		\end{equation}
		where we have written $\beta=(x-T)(Q)/g(P_T)$. Note that $\beta^2\in L$, by the above equation. As we have noted, we then have that
		\begin{equation*}
			\kappa\circ w\circ \delta(P)=\kappa(\beta^\sigma/\beta)\equiv\beta^2\pmod{(L^*)^2}.
		\end{equation*}
		Finally, we compute using the definitions:
		\begin{align*}
			\beta^2&=\frac{(x-T)(2Q)}{g(2P_T)}\\
			&=\frac{(x-T)(2Q)}{(x-T)(2Q-P)}\quad(\text{By Weil reciprocity})\\
			&=(x-T)(P),
		\end{align*}
		so $\kappa\circ w\circ\delta(p)\equiv (x-T)(P)\pmod{(L^*)^2}$, as we wanted.
	\end{proof}
	Combining \thmref{kwisomorphism} with \thmref{x-Tinjection}, we can finally obtain a description of the Selmer group $S^2(J/K)$ which is easier to compute. Using the isomorphism $\kappa\circ w$, we can replace the term $H^1(G_K,J[2])$ in \eqref{seldiag} by the kernel of the norm in $L^*/(L^*)^2$. Then we replace the composition $\kappa\circ w\circ\delta$ by $x-T$. The exact same reasoning applies to $K_\P$ for all $\P$, so if we denote by $L_\P$ the algebra $K_\P[T](f(T))$ we obtain a commutative diagram:
	\begin{equation*}
		\begin{tikzcd}
			0\arrow[r]
			&J(K)/2J(K)\arrow{r}{x-T}\arrow{d}
			&L^*/(L^*)^2\arrow{d}{\prod\text{res}_\P}
			\\
			0\arrow[r]
			&\displaystyle \prod_\P J(K_\P)/2J(K_\P)\arrow{r}{x-T}
			&\displaystyle \prod_\P L_\P^*/(L_\P^*)^2
		\end{tikzcd}
	\end{equation*}
	Tracing the definitions of the maps involved, we find that the Selmer group $S^2(J/K)$ is (isomorphic to) the subgroup of the kernel of the norm from $L^*/(L^*)^2$ to $K^*/(K^*)^2$ whose restriction to each $L_\P^*/(L_P^*)^2$ lies in the image of $x-T$. Moreover, we can study the image of $H^1_S(G_K,J[2])$ through the isomorphism $\kappa\circ w$. The algebra $L$ can be written as a product of fields $L_i$ that are finite extensions of $K$. The number of fields in this product is at most $d$, since $f(T)$ may have nonlinear irreducible factors over $K[T]$. The image of $H^1_S(G_K,J[2])$ is then the group of elements in the kernel of the norm of $L^*/(L^*)^2$ with the following property: if $\lambda\in L$ and $\lambda_i$ is its image in the field $L_i$, the extension $L_i(\sqrt{\lambda_i})$ is ramified only at the primes of $S$. Since we know that the Selmer group lies inside $H^1_S(G_K,J[2])$, the elements of the Selmer group also have this property. We will use these tools in the next chapter to compute the Selmer group of some Jacobian varieties.
	
	
	
	
	
	
	
	
	
	
	
	
	
	
	
	
	
	
	\newpage
	\section{Families}
	
	In this chapter, we will use the theoretical tools that we have developed in the previous chapters to give some bounds on the rank of a family of hyperelliptic curves.
	
	
	Consider the following situation: $r,s$ are distinct odd prime numbers, so we assume without loss of generality that $r>s$, and such that $r+s=2^a\cdot p$ and $r-s=2^b\cdot q$, where $a,b$ are integers and $p,q$ are either odd primes or 1.
	
	We then consider the family of hyperelliptic curves:
	\begin{equation*}
		C_{r,s}: y^2=x\cdot(x^2-(r+s)^2)\cdot(x^2-(r-s)^2)
	\end{equation*}
	
	Throughout the discussion, we will denote:
	\begin{equation*}
		f_{r,s}(x)=x\cdot(x^2-(r+s)^2)\cdot(x^2-(r-s)^2),
	\end{equation*}
	
	so that $C_{r,s}:y^2=f_{r,s}$. We will study the group $J_{r,s}(\Q)$, where $J_{r,s}$ is the Jacobian of the curve $C_{r,s}$. Specifically, we will give bounds for the Mordell-Weil of $J_{r,s}(\Q)$ depending on some congruence relations for $r,s,p,q,a$ and $b$.
	
	To begin with, note that the 2-torsion of $J_{r,s}(\Q)$ is isomorphic to $(\Zmod{2})^4$, with a basis given by $(0,0)-\infty$, $(r+s,0)-\infty$, $(-r-s,0)-\infty$ and $(r-s,0)-\infty$. Observe that the 2-torsion point corresponding to the fifth root of $f_{r,s}$ is a linear combination of the other given 2-torsion points, since:
	\begin{equation*}
		\Div(y)=\sum_{i=1}^{5}((\alpha_i,0)-\infty),
	\end{equation*}
	where the $\alpha_i$ are the distinct roots of $f_{r,s}$.
	
	
	\textcolor{red}{COMENTARI SOBRE LA TORSIÓ}
	
	
	
	By the Mordell-Weil theorem, \textcolor{red}{FALTA CITAR} we have
	\begin{equation*}
		J_{r,s}(\Q)\cong \Z^r\oplus \left( \Zmod{2}\right)^4,
	\end{equation*}
	for some $r\ge 0$. Therefore, we have that:
	\begin{equation*}
		2J_{r,s}(\Q)\cong (2\Z)^r,
	\end{equation*}
	so we find that:
	\begin{equation*}
		J_{r,s}/2J_{r,s}\cong (\Zmod{2})^r\oplus(\Zmod{2})^4\cong (\Zmod{2})^{r+4}.
	\end{equation*}
	For this reason, in order to calculate the Mordell-Weil rank of $J_{r,s}$ we have to compute 2-dimension of the group $J_{r,s}/2J_{r,s}$.
	
	We will first explain some considerations on the parameters $r,s,p,q,a$ and $b$. Since $r$ and $s$ are odd and $r>s$, the values $r+s$ and $r-s$ are even and strictly positive. Since $p$ and $q$ are odd, this forces that $a,b\ge 1$. Next we will show that except for a particular case, the values $r,s,p,q$ are different.
	
	If $p$ is equal to $q$, then we have:
	\begin{equation*}
		2r=(r+s)+(r-s)=(2^a+2^b)p.
	\end{equation*}
	Since $a,b\ge 1$, this implies that $r=(2^{a-1}+2^{b-1})p$. Since $2^{a-1}+2^{b-1}\ge 2$, if we have $p>1$ then $r$ is not a prime, so $p=1$, and thus $q=1$. In this case, if $a=b$ we have $r=2\cdot 2^{a-1}=2^a$, which is only prime if $a=1$. But then $r+s=2=r$, which is impossible. Therefore we must have $a\neq b$, and since $r+s>r-s$, we must have $a>b$. If $b>1$, then $r=2^{a-1}+2^{b-1}>1+1=2$ and it is even, so it cannot be prime. Thus we must have $b=1$ and $a>1$, so $r=2^{a-1}+1$ and $r-s=2$, so $s=2^{a-1}-1$. Observe that then $s$ is a Mersenne prime, so we know that $a-1$ is prime. If $a-1=2$, so $a=3$, we have the specific case:
	\begin{equation*}
		\fbox{C1}: r=5,s=3,p=1,q=1,
	\end{equation*}
	which will be dealt later. Now if $a-1$ is odd, we have:
	\begin{gather*}
		s=2^{a-1}-1\equiv -1\equiv 1\pmod 2,\\
		s=2^{a-1}-1\equiv (-1)^{a-1}-1\equiv -2\equiv 1\pmod 3,
	\end{gather*}
	so $s\equiv 1\pmod 6$. Therefore, $r=s+2\equiv 3\pmod6$, which is only prime if $r=3$ but then $s=1$ which is impossible. In conclusion, except for the case $C1$, $p$ and $q$ are different. Now we rule out every other possibility:
	
	Assume that $p$ is equal to $r$. Then $2r>r+s=2^ar$, so $1>2^{a-1}\ge 1$, which is impossible.
	
	Assume that $p$ is equal to $s$. Then $r+s=2^as$, so $r=s(2^a-1)$. If $a>1$, then $2^a-1>1$, so $r$ is not prime. If $a=1$, then $r=s$ which is also impossible.
	
	Assume that $q$ is equal to $s$. Then $r-s=2^bs$, so $r=(2^b+1)s$. Note that $s>1$ and, for every $b\ge 1$, $2^b+1>1$, so $r$ cannot be prime, which is impossible.
	
	Finally, assume that $q$ is equal to $r$. Then $2^br=r-s<r<2^br$, which is impossible.
	
	
	Now our choice for the pairs $r$ and $s$ will become clear: we want to determine the prime of bad reduction for $J_{r,s}$, which correspond to the primes dividing the discriminant of $f_{r,s}(x)$, so we compute it:
	\begin{equation*}
		\Delta_{r,s}=(r+s)^6\cdot(r-s)^6\cdot s^4\cdot r^4\cdot 2^{12}.
	\end{equation*}
	
	From our initial assumption on the values of $r+s$ and $r-s$, we find that the only primes dividing $\Delta_{r,s}$ are $r,s,p,q$ and $2$. From our previous discussion, except for a specific case these five numbers are distinct from each other. Therefore our family of hyperelliptic curves has a uniform bound on the number of primes of bad reduction, which will be useful for our computations.
	
	We now will deal with the specific case $C1$: $r=5$, $s=3$, so that the curve becomes:
	\begin{equation*}
		C_{5,3}: y^2=x\cdot(x^2-8^2)\cdot(x^2-2^2).
	\end{equation*}
	Since $p=q=1$, the primes of bad reduction for $J_{r,s}$ are just $2,3$ and $5$. In the notation of the previous chapter, we take $S=\{\infty,2,3,5\}$ and $H=\langle -1,2,3,5\rangle$, which is isomorphic to $\GF2^4$. We will work in the vector space $H^5\cong \GF2^{20}$, specifically in the subgroup $G$ which is the kernel of the norm function. Since we are working with a specific equation, we can find an extra point of the curve: when $x=-4$, the equation becomes $y^2=2304$, which has two solutions $y=\pm 48$ in $\Q$. We will be able to see that $J_{5,3}$ has rank 1 and that the point $(-4,48)-\infty$ is a generator.
	
	We now compute the image of the known rational points inside $H^5$, and present the results in a table. Since we have the natural decomposition:
	\begin{equation*}
		L=\Q[T]/(f_{r,s}(T))\cong \Q^5
	\end{equation*}
	which descends to a decomposition
	\begin{equation*}
		L^*/(L^*)^2\cong \Q^*/(\Q^*)^2\oplus \overset{(5}{\ldots}\oplus \Q^*/(\Q^*)^2,
	\end{equation*}
	we can compute the map $X-T$ from Chapter \textcolor{red}{NUMERO DEL CAPITOL ANTERIOR}
	
	
	Rational points
	
	
	{
		\centering
		\begin{tabular}{c|c|c|c|c|c}
			 & $X+8$ & $X+2$ & $X-0$ & $X-2$ & $X-8$\\ \hline
			$(-8,0)-\infty$&$30$ & $-6$ & $-2$ & $-10$ & $-1$\\ \hline
			$(-2,0)-\infty$&$6$ & $-30$ & $-2$ & $-1$ & $-10$\\ \hline
			$(0,0)-\infty$&$2$ & $2$ & $1$ & $-2$ & $-2$\\ \hline
			$(2,0)-\infty$&$10$ & $1$ & $2$ & $-30$ & $-6$\\ \hline
			$(-4,48)-\infty$&$1$ & $-2$ & $-1$ & $-6$ & $-3$
		\end{tabular}\par
	}
	
	Real points
	
	
	{
	\centering
	\begin{tabular}{c|c|c|c|c|c}
		& $X+8$ & $X+2$ & $X-0$ & $X-2$ & $X-8$\\ \hline
		$(-2,0)-\infty$&$1$ & $-1$ & $-1$ & $-1$ & $-1$\\ \hline
		$(0,0)-\infty$&$1$ & $1$ & $1$ & $-1$ & $-1$\\ \hline
	\end{tabular}\par
}	
	
	$\Q_2$ points
	
	
{
	\centering
	\begin{tabular}{c|c|c|c|c|c}
		& $X+8$ & $X+2$ & $X-0$ & $X-2$ & $X-8$\\ \hline
		$(-8,0)-\infty$&$-2$ & $-6$ & $-2$ & $6$ & $-1$\\ \hline
		$(-2,0)-\infty$&$6$ & $2$ & $-2$ & $-1$ & $6$\\ \hline
		$(0,0)-\infty$&$2$ & $2$ & $1$ & $-2$ & $-2$\\ \hline
		$(2,0)-\infty$&$-6$ & $1$ & $2$ & $2$ & $-6$\\ \hline
		$(-4,48)-\infty$&$1$ & $-2$ & $-1$ & $-6$ & $-3$\\ \hline
		$(-3,\varepsilon)-\infty$ &$-3$ & $-1$ & $-3$ & $3$ & $-3$
	\end{tabular}\par
}

	$\Q_3$ points

{
	\centering
	\begin{tabular}{c|c|c|c|c|c}
		& $X+8$ & $X+2$ & $X-0$ & $X-2$ & $X-8$\\ \hline
		$(-8,0)-\infty$&$3$ & $3$ & $1$ & $-1$ & $-1$\\ \hline
		$(-2,0)-\infty$&$-3$ & $-3$ & $1$ & $-1$ & $-1$\\ \hline
		$(0,0)-\infty$&$-1$ & $-1$ & $1$ & $1$ & $1$\\ \hline
		$(2,0)-\infty$&$1$ & $1$ & $-1$ & $-3$ & $3$\\ \hline
		$(-4,48)-\infty$&$1$ & $1$ & $-1$ & $3$ & $-3$
	\end{tabular}\par
}

	$\Q_5$ points

{
	\centering
	\begin{tabular}{c|c|c|c|c|c}
		& $X+8$ & $X+2$ & $X-0$ & $X-2$ & $X-8$\\ \hline
		$(-8,0)-\infty$&$5$ & $1$ & $2$ & $10$ & $1$\\ \hline
		$(-2,0)-\infty$&$1$ & $5$ & $2$ & $1$ & $10$\\ \hline
		$(0,0)-\infty$&$2$ & $2$ & $1$ & $2$ & $2$\\ \hline
		$(2,0)-\infty$&$10$ & $1$ & $2$ & $5$ & $1$\\ \hline
		$(-4,48)-\infty$&$1$ & $2$ & $1$ & $1$ & $2$
	\end{tabular}\par
}

Basis for $\Q$ points that begin with 1:

{
	\centering
	\begin{tabular}{c c c c c}
		$1$ & $10$ & $2$ & $-6$ & $30$\\ \hline
		$1$ & $-2$ & $-1$ & $6$ & $-3$
	\end{tabular}\par
}

Basis for $\Q_3$ points that begin with 1:

{
	\centering
	\begin{tabular}{c c c c c}
		$1$ & $1$ & $-1$ & $-3$ & $3$\\ \hline
		$1$ & $1$ & $-1$ & $3$ & $-3$
	\end{tabular}\par
}

note that second entry is 1 in $\Q_3^*/(\Q_3^*)^2$, so second entry is contained in $\langle -2,10\rangle$. In the basis of the known points that begin with 1, we have all the possibilities so we must show that there is only $2^{2-2}=1$ element that begins with $1,1$, namely the identity. Also note that the third entry is in $\langle -1\rangle$ of $\Q_3$, so it is in $\langle -1,2,5\rangle$.

Basis for $\Q_2$ points that begin with 1,1:

{
	\centering
	\begin{tabular}{c c c c c}
		$1$ & $1$ & $6$ & $2$ & $3$
	\end{tabular}\par
}

so third entry is contained in $\langle 6\rangle$ of $\Q_2$, and thus third entry in $\Q$ is contained in $\langle 6,-15\rangle$.

so fourth entry is contained in $\langle 2\rangle$ of $\Q_2$, and thus fourth entry in $\Q$ is contained in $\langle 2,-15\rangle$.

so fifth entry is contained in $\langle 3\rangle$ of $\Q_2$, and thus fifth entry in $\Q$ is contained in $\langle 3,-15\rangle$.



Basis for $\Q_5$ points that begin with 1,1: just the identity, so a points that begins with $1,1$ has the other three entries in $\langle -1,6\rangle$.

Comparing with $\Q_2$, we see that fourth entry is in $\langle 2,-15\rangle\cap \langle -1,6\rangle= 1$ and fifth entry is in $\langle 3,-15\rangle\cap \langle -1,6\rangle= 1$ so fourth and fifth are 1. Third entry is in $\langle 6,-15\rangle\cap \langle 6,-1\rangle=\langle 6\rangle$, but also in $\langle -1,2,5\rangle$ from $\Q_3$, so it is one, and we are done.

	\newpage
	\appendix
	\section{Abelian Varieties}
	In this appendix, we will make a brief survey on the definitions and fundamental results for abelian varieties. Some of the theorems presented here will be used in Chapters 1 and 2. We will follow \cite{Cornell_Silverman_Artin_1986}, \cite{Hartshorne_1977}, and \cite{Hindry_Silverman_2000}.

Throughout the exposition, $K$ will denote a field of characteristic zero, and $\overline{K}$ will denote its algebraic closure.
\begin{definition}
	A \emph{group variety} defined over $K$ is a variety $V$ defined over $K$, together with a point $\varepsilon\in V(K)$ and morphisms $\mu:V\times V\rightarrow V$, $i:V\rightarrow V$ such that the following diagrams
	\begin{equation*}
		\begin{tikzcd}
			V\times V\times V\arrow{r}{(\mu,\text{Id})}\arrow{d}[left]{(\text{Id},\mu)}
			&V\times V\arrow{d}{\mu}\\
			V\times V\arrow{r}{\mu}
			& V
		\end{tikzcd}\text{(Associativity)},
	\end{equation*}
	\vspace{20pt}
	\begin{equation*}
		\begin{tikzcd}
				V\arrow{r}{(\text{Id},e)}\arrow[bend right]{dr}[below]{\text{Id}}
				&V\times V\arrow{d}{\mu}
				&V\arrow{l}[above]{(e,\text{Id})}\arrow[bend left]{dl}{\text{Id}}\\
				&V
		\end{tikzcd}\text{(Identity)},
	\end{equation*}
	and
		\vspace{15pt}
	\begin{equation*}
		\begin{tikzcd}
			V\arrow{r}{(\text{Id},i)}\arrow[bend right]{dr}[below]{e}
			&V\times V\arrow{d}{\mu}
			&V\arrow{l}[above]{(i,\text{Id})}\arrow[bend left]{dl}{e}\\
			&V
		\end{tikzcd}\text{(Inverse)}
	\end{equation*}
	commute, where $e:V\rightarrow V$ is the constant map that sends every point $P\in V$ to $\varepsilon$. This means that the pair $(V,\mu)$ is a group.
\end{definition}
In a more abstract language, we say that a group variety over $K$ is a group object in the category of varieties over $K$.

\begin{example}
	The affine line $\mathbb{A}^1(K)$, together with the usual sum on $K$, is a group variety.
\end{example}

If $V$ is a group variety over $K$, we can define, for any point $P\in V(K)$, the \emph{left (resp. right) translate by $P$} as the maps:
\begin{alignat*}{2}
	L_P:V&\rightarrow V \qquad\qquad R_P:V&&\rightarrow V\\
	Q&\mapsto \mu(P,Q),\qquad\qquad Q&&\mapsto \mu(Q,P),
\end{alignat*}
which are defined over $K$. Using the definitions, we can see that the translation maps are isomorphisms of varieties (with the inverse given by $L_{i(P)}$ and $R_{i(P)}$, respectively). From this, we can see that any group variety is smooth: any variety has some non-singular point $P$ (in fact it has an open subset consisting of non-singular points). If $Q$ were a singular point in $V$, then the point $L_{\mu(P,i(Q))}(Q)$ would also be singular, but
\begin{equation*}
	L_{\mu(P,i(Q))}(Q)=\mu(\mu(P,i(Q)),Q)=\mu(P,\mu(i(Q),Q))=\mu(P,\varepsilon)=P.
\end{equation*}

\begin{lemma}[Rigidity lemma]\label{rigidity}
	Let $X$ be a projective variety, let $Y$ and $Z$ be any varieties, and let $f:X\times Y\rightarrow Z$ be a morphism. Suppose that there is some point $y_0\in Y$ such that $f$ is constant on $X\times\{y_0\}$. Then, $f$ is constant on every set of the form $X\times\{y\}$, $y\in Y$. Furthermore, if $f$ is also constant on some $\{x_0\}\times Y$, then it is constant on all $X\times Y$.
\end{lemma}
\begin{proof}
	The proof relies on the crucial fact that projective varieties are complete (see \cite{Hartshorne_1977}). Consider the projection $p:X\times Y\rightarrow Y$, which is closed. Let $z_0$ be the constant value of $f$ on the set $X\times\{y_0\}$. Let $U$ be an affine open neighborhood of $z_0$ in $Z$. Therefore, $Z\setminus U$ is closed, so $f^{-1}(Z\setminus U)$ is closed. Hence, the set $W=p(f^{-1}(Z\setminus U))$ is closed in $Y$. Since $y_0\not\in W$ (otherwise some $x\in X$ would satisfy $f(x,y_0)\neq z_0$), we have that $Y\setminus W$ is a nonempty open subset of $Y$, so it is dense. Now if $y$ is any point in $Y\setminus W$, we have that $f(x,y)$ is in $U$ for all $x\in X$, so $f(X\times\{y\})$ is contained in $U$. Since $X\times\{y\}\cong X$ is complete and $U$ is affine, we must have that $f(X\times\{y\})$ is constant. As $Y\setminus W$ is dense, the same holds for any $y\in Y$, so we are done. The second assertion follows clearly from the first one.
\end{proof}
\begin{definition}
	An \emph{abelian variety} $A$ is a projective group variety.
\end{definition}
\begin{example}
	An elliptic curve $E$, together with its usual group structure, is an abelian variety.
\end{example}
Note that in our definition of group variety we have not assumed the group law to be commutative. Using the rigidity lemma and the fact that abelian varieties are projective, we will be able to show that the group of an abelian variety is an abelian group. For now, we write the group law multiplicatively.
\begin{proposition}\label{morphism}
	Let $A,B$ be abelian varieties, and let $\phi:A\rightarrow B$ be a morphism. Then $\phi$ is the composition of a translation and an homomorphism.
\end{proposition}
\begin{proof}
	Denote by $\varepsilon_A$ and $\varepsilon_B$ the identity element of $A$ and $B$, respectively. After a possible translation, we may assume that $\phi(\varepsilon_A)=\varepsilon_B$. Consider the function
	\begin{align*}
		f:A\times A&\rightarrow B\\
		(x,y)&\mapsto \phi(xy)\phi(y)^{-1}\phi(x)^{-1}.
	\end{align*}
	Observe that $f(x,\varepsilon_A)=\phi(x)\phi(\varepsilon_A)^{-1}\phi(x)^{-1}=\varepsilon_B$ for all $x\in A$. Similarly, $f(\varepsilon_A,y)=\phi(y)\phi(y)^{-1}\phi(\varepsilon_A)^{-1}=\varepsilon_B$ for all $y\in A$, so \lemref{rigidity} implies that $f$ is constant (with constant value $\varepsilon_B$). Hence, we have for all $x,y\in A$:
	\begin{equation*}
		\varepsilon_B=\phi(xy)\phi(y)^{-1}\phi(x)^{-1},
	\end{equation*}
	so
	\begin{equation*}
		\phi(xy)=\phi(x)\phi(y),
	\end{equation*}
	i.e., $\phi$ is a homomorphism.
\end{proof}

\begin{corollary}
	The group law of an abelian variety $A$ is commutative.
\end{corollary}
\begin{proof}
	By assumption, the inversion map $i:A\rightarrow A$ is a morphism. By \propref{morphism}, it is the composition of a translation and a homomorphism, say $i=\varphi\circ\tau$. Since $i(\varepsilon_A)=\varepsilon_A$, the translation $\tau$ is the identity. Therefore, $i$ is a homomorphism, which is equivalent to $A$ being an abelian group.
\end{proof}
From this point on, we will write the group law of an abelian variety additively. The right notion of morphism of abelian varieties is that of an isogeny.
\begin{definition}
	Let $A$, $B$ be abelian varieties of the same dimension. A morphism of varieties $\phi\in\Hom(A,B)$ is an \emph{isogeny} if it maps the identity element of $A$ to the identity element of $B$, and is either constant or surjective with finite kernel. The degree of $\phi$ is defined to be the cardinality of $\ker(\phi)$. (This definition of degree is not the general one, but in characteristic zero it coincides with the usual definition)
\end{definition}
Note that the fact that isogenies preserve the identity, together with \propref{morphism}, implies that isogenies are group morphisms. If an isogeny $\phi:A\rightarrow B$ is constant, then it must be the constant map to the identity of $B$. If it is surjective, observe that there exists a nonempty open subset $U$ of $B$ such that $\dim{\phi^{-1}}(\{p\})=\dim A-\dim B=0$, for all $p\in U$, by the theorem on the dimension of the fibers. Choose any $p\in U$, and find a point $a_0\in A$ such that $\phi(a_0)=p$ (note that $\phi$ is surjective). Then, we have the equality of sets
\begin{equation*}
	\phi^{-1}(\{p\})-a_0=\{a-a_0\in A:\phi(a)=p\}=\{a\in A:\phi(a)=0\}=\ker\phi,
\end{equation*}
because $\phi$ is a group morphism. Since translate maps are bijective, we find that $\dim(\ker\phi)=\dim(\phi^{-1}(\{p\}))=0$. Zero dimensional subvarieties are finite sets of points, so we find that the kernel is finite (so we do not need that condition on the definition of isogeny).

For any abelian variety $A$ and any integer $n$, we can consider the multiplication by $n$ map which we denote by $[n]$. We will study these maps later, but they constitute the most important examples of isogenies of an abelian variety.

We want to study divisors on abelian varieties, in a way to obtain relations among them that exhibit the group structure of $A$. The fundamental result in this line, which once more reflects the rigidity of abelian varieties, is the theorem of the cube.

\begin{theorem}[Theorem of the cube]\label{thmcube}
	Let $X$, $Y$ and $Z$ be projective varieties, and let $(x_0,y_0,z_0)\in X\times Y\times Z$. Let $D$ be a divisor that is linearly equivalent to the trivial class on the sides:
	\begin{equation*}
		X\times Y\times\{z_0\},\quad X\times\{y_0\}\times Z,\quad\text{and } \{x_0\}\times Y\times Z.
	\end{equation*}
	Then $D$ is linearly equivalent to the trivial class on $X\times Y\times Z$.
\end{theorem}
\begin{proof}
	See \cite{Cornell_Silverman_Artin_1986}, Chapter 5 or \cite{Mumford_1970}.
\end{proof}
In the case of abelian varieties, the theorem of the cube has a more particular expression. We fix some useful notation. For any nonempty subset $I\subset\{1,2,3\}$, we define a map
\begin{equation*}
	s_I:A\times A\times A\rightarrow A,\quad s_I(x_1,x_2,x_3)=\sum_{i\in I}x_i.
\end{equation*}
With this notation, we have the following theorem.
\begin{theorem}[Theorem of the cube for abelian varieties]\label{thmcubeabelian}
	Let $A$ be an abelian variety, and let $D\in\Div(A)$ be a divisor. Then
	\begin{equation*}
		s_{123}^*D-s_{12}^*D-s_{13}^*D-s_{23}^*D+s_1^*D+s_2^*D+s_3^*D\sim 0.
	\end{equation*}
\end{theorem}
\begin{proof}
	Define $\text{cube}(D)$ as the left hand side of the expression in the theorem, and note that we have
	\begin{equation*}
		\text{cube}(D)=-\sum_{I\subset\{1,2,3\}}(-1)^{|I|}s_I^*(D).
	\end{equation*}
	We want to apply \thmref{thmcube} to the divisor $\text{cube}(D)\in\Div(A\times A\times A)$. By symmetry, it suffices to show that it becomes trivial when restricted to some $A\times A\times\{z_0\}$, so we choose $z_0=0$. Therefore, consider the injection $i:A\times A\rightarrow A\times A\times A$ given by $i(x,y)=(x,y,0)$, so we have to show that $i^*(\text{cube}(D))\sim 0$. We have the following equalities:
	\begin{align*}
		s_{123}\circ i(x,y)&=x+y=s_{12}\circ i(x,y),\\
		s_{23}\circ i(x,y)&=y=s_2\circ i(x,y),\\
		s_{13}\circ i(x,y)&=x=s_1\circ i(x,y), \text{ and}\\
		s_3\circ i(x,y)&=0.
	\end{align*}
	Therefore, observe that
	\begin{equation*}
		i^*(\text{cube}(D))=-\sum_{I\subset\{1,2,3\}}(-1)^{|I|}(s_I\circ i)^*(D)\sim 0,
	\end{equation*}
	because the pairs corresponding to $(s_{123},s_{12})$, $(s_{23},s_2)$, and $(s_{13},s_1)$ cancel each other out, and the term corresponding to $s_3$ is zero,	so we get the result.
\end{proof}
We will also need an important result, which is a version of the theorem of the cube for two varieties.
\begin{theorem}[Seesaw principle]
	Let $X$ and $Y$ be two varieties, and let $D\in\Div(X\times Y)$ be a divisor. For $x_0\in X,y_0\in Y$, define the maps $i_{x_0}:Y\rightarrow X\times Y$ and $i_{y_0}:X\rightarrow X\times Y$ as $i_{x_0}(y)=(x_0,y)$ and $i_{y_0}(x)=(x,y_0)$. Suppose that the following conditions hold:
	\begin{itemize}
		\item[(a)] $D$ becomes linearly equivalent to zero when restricted to $\{x\}\times Y$ for each $x\in X$, that is, the pullback $i_x^*D$ is linearly equivalent to zero for all $x\in X$, and
		\item[(b)] $D$ becomes linearly equivalent to zero when restricted to $X\times\{y_0\}$, for some $y_0\in Y$, that is, the pullback $i_{y_0}^*D$ is linearly equivalent to zero for some $y_0\in Y$.
	\end{itemize}
	Then, $D$ is linearly equivalent to zero.
\end{theorem}
For proofs of the general theorem of the cube and the seesaw principle, we refer to \cite{Mumford_1970}.
Now we can deduce some relations between divisors of an abelian variety.
\begin{corollary}\label{threemorphisms}
	Let $A$ be an abelian variety, $V$ an arbitrary variety, and $f,g,h:V\rightarrow A$ three morphisms. Then for any divisor $D\in\Div(A)$, we have:
	\begin{equation*}
		(f+g+h)^*D-(f+g)^*D-(f+h)^*D-(h+g)^*D+f^*D+g^*D+h^*D\sim 0
	\end{equation*}
	in $\Div(V)$.
\end{corollary}
Note that the above relation is in $V$, and the sum of maps is considered pointwise.
\begin{proof}
	Let $(f,g,h):V\rightarrow A\times A\times A$ be the map $(f,g,h)(x)=(f(x),g(x),h(x))$. In the notation of the theorem of the cube for abelian varieties, observe that:
	\begin{gather*}
		(f,g,h)^*(\text{cube}(D))=-\sum_{I\subset\{1,2,3\}}(-1)^{|I|}(s_I\circ(f,g,h))^*(D)\\
		=(f+g+h)^*D-(f+g)^*D-(f+h)^*D-(h+g)^*D+f^*D+g^*D+h^*D.
	\end{gather*}
	But $\text{cube}(D)\sim 0$ by \thmref{thmcubeabelian}, so the result follows.
\end{proof}

\begin{corollary}\label{mumfordformula}
	Let $D$ be a divisor on an abelian variety $A$, and let $[m]:A\rightarrow A$ be the multiplication by $m$ map. Then, we have
	\begin{equation*}
		[m]^*D\sim\frac{m^2+m}{2}D+\frac{m^2-m}{2}[-1]^*D.
	\end{equation*}
	In particular, if $D$ is symmetric, we have $[m]^*D\sim m^2D$.
\end{corollary}
\begin{proof}
	Observe that the formula is trivial for $m=-1,0$ and $1$. Now, use \corref{threemorphisms} with $f=[m]$, $g=[1]$ and $h=[-1]$ to obtain:
	\begin{align}\label{mumfordinduction}
		[m]^*D-[m+1]^*D-[m-1]^*D+[m]^*D+D+[-1]^*D\\
		=2[m]^*D-[m+1]^*D-[m-1]*D+D+[-1]^*D\sim 0.
	\end{align}
	Now we can do induction on $m$: assume that we know the result for $m\ge1$ and $m-1$. Observe the following identities:
	\begin{align*}
		\frac{(m-1)^2+(m-1)}{2}&=\frac{m^2-m}{2},\\
		\frac{(m-1)^2-(m-1)}{2}&=\frac{m^2-3m+2}{2},\\
		(m^2+m)-\frac{m^2-m}{2}+1&=\frac{(m^2+3m+2)}{2}=\frac{(m+1)^2+(m+1)}{2}, \text{ and}\\
		(m^2-m)-\frac{m^2-3m+2}{2}+1&=\frac{m^2+m}{2}=\frac{(m+1)^2-(m+1)}{2}.
	\end{align*}
	
	Now use these identities and the relation \eqref{mumfordinduction} to compute:
	\begin{align*}
		&[m+1]^*D\sim2[m]^*D-[m+1]^*D+D+[-1]^*D\\
		&\sim(m^2+m)D+(m^2-m)[-1]^*D-\frac{m^2-m}{2}D-\frac{m^2-3m+2}{2}[-1]^*D+D+[-1]^*D\\
		&\sim\frac{(m+1)^2+(m+1)}{2}D+\frac{(m+1)^2-(m+1)}{2}[-1]^*D,
	\end{align*}
	as we wanted. Induction for negative integers is analogous, and the assertion on symmetric divisors now follows trivially.
\end{proof}

Additionally, we can deduce an important fact about translation maps $t_a$ for any $a\in A$ (note that since abelian varieties are abelian group, we do not need to make a distinction between right and left translates). This theorem can also be proved independently and then used to prove the theorem of the cube.
\begin{theorem}[Theorem of the square]\label{thmsquare}
	Let $A$ be an abelian variety and consider, for any $a\in A$, the translate map $t_a:A\rightarrow A$. Then for any divisor $D\in\Div(A)$, we have:
	\begin{equation*}
		t_{a+b}^*(D)+D\sim t_a^*D+t_b^*D\text{ for all }a,b\in A.
	\end{equation*}
\end{theorem}
\begin{proof}
	Consider the morphisms $f,g,h:A\rightarrow A$ given by $f(x)=x$, $g(x)=a$ and $h(x)=b$. Then, \corref{threemorphisms} yields:
	\begin{align*}
		(x+a+b)^*D&-(x+a)^*D-(x+b)^*D-(a+b)^*D+x^*D+a^*D+b^*D\\
		&\sim (t_{a+b})^*D-t_a^*D-t_b^*D+D\sim 0,
	\end{align*}
	which is what we wanted.
\end{proof}

As we claimed earlier, the maps $[m]$ are isogenies. We study them in the next theorem.
\begin{theorem}\label{isogenymult}
	Let $A$ be an abelian variety of dimension $g$ over an algebraically closed $K$. Then the multiplication by $m$ map $[m]:A\rightarrow A$ is an isogeny of degree $n^{2g}$, and
	\begin{equation*}
		A[m]=\ker[m]\cong\left(\Zmod{m}\right)^{2g}.
	\end{equation*}
\end{theorem}
\begin{proof}
	The multiplication map $\mu:A\times A\rightarrow A$ induces a linear map on the tangent space of $A$ at $(0,0)$, $\mu_*:T_0(A)\times T_0(A)\rightarrow T_0(A)$, which corresponds to addition on the vector space $T_0(A)$. Therefore the map $[m]:A\rightarrow A$ induces the multiplication by $m$ on $T_0(A)$. Since we are working in characteristic zero, this map is an isomorphism. Thus, we have that $\dim(A)=\dim([m]A)$, so $A=[m]A$, i.e., it is surjective. Since $[m]0=0$, we conclude that $[m]$ is an isogeny. The degree of the maps $[m]$ and the structure of the $m$-torsion can be studied in $\mathbb{C}$ from the theory of complex tori, and then use the Lefschetz principle to obtain the result in any algebraically closed field of characteristic zero. Instead, we give an independent proof here. We will use the following fact on intersection indices of divisors: if $A$ is an abelian variety of dimension $g$ over an algebraically closed field $\overline{K}$, and $\phi:A\rightarrow A$ is an isogeny, then
	\begin{equation*}
		(\phi^*D_1,\ldots,\phi^*D_g)_A=\deg(\phi)(D_1,\ldots,D_g)_A,
	\end{equation*}
	for any divisors $D_1,\ldots,D_g\in\Div(A)$, and if $D$ is an ample divisor, then $D^g=(D,\ldots,D)$ is effective. For a proof, see \cite{Hindry_Silverman_2000}.
	To use this fact, take any symmetric ample divisor on $A$, and compute:
	\begin{equation*}
		\deg([m])D^g=([m]^*D)^g=(m^2D)^g=m^{2g}D^g.
	\end{equation*}
	Since $D^g>0$, we deduce that $\deg([m])=m^{2g}$, as we wanted. In characteristic zero, we have then that $|\ker([m])|=m^{2g}$. But the same argument works for any integer, in particular for those dividing $m$. Using elementary group theory, it can be seen that a finite abelian group whose order is a perfect power $m^r$, and with $n$ torsion of order $n^r$, for any integer $n$ dividing $m$, must be isomorphic to $(\Zmod{m})^r$. In our case we find that $\ker([m])\cong(\Zmod{m})^{2g}$, as we wanted.
\end{proof}
In abelian varieties, symmetric divisors can be characterized by the following result.
\begin{proposition}\label{abeliansymmetric}
	Let $A$ be an abelian variety, and let $D\in\Div(A)$ be a divisor. Then $D$ is symmetric if, and only if,
	\begin{equation}\label{parallelogramequation}
		s_{12}^*D+d_{12}^*D\sim2p_1^*D+2p_2^*D,
	\end{equation}
	where $s_{12}:A\times A\rightarrow A$ is $s_{12}(x,y)=x+y$, $d_{12}:A\times A\rightarrow A$ is $d_{12}(x,y)=x-y$ and the $p_i:A\times A\rightarrow A$ are the natural projections.
\end{proposition}
\begin{proof}
	Let $j:A\rightarrow A\times A$ be the map $j(x)=(0,x)$. We have:
	\begin{align*}
		s_{12}\circ j(x)&=x,\\
		d_{12}\circ j(x)&=-x,\\
		p_1\circ j(x)&=0,\text{ and}\\
		p_2\circ j(x)&=x.
	\end{align*}
	Therefore, we have:
	\begin{align*}
		j^*(s_{12}^*D+d_{12}^*D-2p_1^*D-2p_2^*D)&=(s_{12}\circ j)^*D+(d_{12}\circ j)^*D-2(p_1\circ j)^*D-2(p_2\circ j)^*D\\
		&=D+[-1]^*D-2D^*\\
		&=[-1]^*D-D.
	\end{align*}
	From this equality, we see that if \eqref{parallelogramequation} holds, then $[-1]^*D-D\sim 0$, so $D$ is symmetric.
	
	Conversely, assume that $D$ is symmetric. For any fixed $a\in A$, consider the map $i:A\rightarrow A\times A$ given by $i_a(x)=(x,a)$. Also consider the translate by $a$ map $t_a$. Then, we have:
	\begin{align*}
		s_{12}\circ i_a(x)=x+a&=t_a(x),\\
		d_{12}\circ i_a(x)=x-a&=t_{-a}(x),\\
		p_1\circ i_a(x)&=x,\text{ and}\\
		p_2\circ i_a(x)&=a.
	\end{align*}
	Using the theorem of the square, we can compute:
	\begin{align*}
		i_a^*(s_{12}^*D+d_{12}^*D-2p_1^*D-2p_2^*D)&=(s_{12}\circ i_a)^*D+(d_{12}\circ i_a)^*D-2(p_1\circ i_a)^*D-2(p_2\circ i_a)^*D\\
		&=t_a^*D+t_{-a}^*D-2D\\
		&\sim t_{a-a}^*D+D-2D\\
		&=2D-2D=0.
	\end{align*}
	We want to apply the seesaw principle to $s_{12}^*D+d_{12}^*D-2p_1^*D-2p_2^*D$ in order to conclude that it is trivial. The above shows that the first property required is fulfilled. Now we want to see that this divisor is trivial on some set $A\times\{y_0\}$, with $y_0\in A$. Choose $0\in A$, and consider the injection $j$ from before. We have to check that $j^*(s_{12}^*D+d_{12}^*D-2p_1^*D-2p_2^*D)$ is trivial. We have seen that it is linearly equivalent to $[-1]^*D-D$, so by our assumption on $D$ being symmetric, we conclude that $s_{12}^*D+d_{12}^*D-2p_1^*D-2p_2^*D$ is trivial when restricted to $A\times\{0\}$, so it is trivial.
\end{proof}
	
	\newpage
	\pagestyle{empty}
	\bibliographystyle{amsplain}
	\bibliography{refs}
	\nocite{*}

\end{document}




